%% ============================================================
%% IEEE TEVC — Paper 2 (revised after the two-stage ILP baseline)
%% Dynamic Kidney Exchange optimization Under Realistic Clinical
%% Constraints for Equitable Allocation
%% ------------------------------------------------------------
%% Revision: the framework is unchanged (volume-ILP-guided NSGA-III).
%% The two-stage ILP is reported as a separate, stronger baseline at
%% maximum volume (modules 12, 13, 16, 17). Numbers come from the
%% module reports. Red \tbd{...} = fill in after rerunning module 17.
%% Figures: upload results/paper/figures_final/*.png to "figures".
%% ============================================================

\documentclass[journal]{IEEEtran}

%% ── PACKAGES ────────────────────────────────────────────────
\usepackage[T1]{fontenc}
\usepackage[utf8]{inputenc}
\usepackage{amsmath, amssymb, amsthm}
\usepackage{algorithm}
\usepackage{algpseudocode}
\usepackage{graphicx}
\usepackage{booktabs}
\usepackage{multirow}
\usepackage{array}
\usepackage{xcolor}
\newcommand{\tbd}[1]{\textcolor{red}{#1}}
\usepackage{url}
\usepackage{cite}
\usepackage{hyperref}
\usepackage{cleveref}      % must be loaded after hyperref
\usepackage{balance}       % balanced columns on the last page

%% ── THEOREM ENVIRONMENTS ────────────────────────────────────
\newtheorem{definition}{Definition}
\newtheorem{remark}{Remark}

%% ── CUSTOM COMMANDS ─────────────────────────────────────────
\newcommand{\pof}{\ensuremath{\mathrm{PoF}_t}}
\newcommand{\gini}{\ensuremath{\mathrm{Gini}}}
\newcommand{\cpra}{\ensuremath{\mathrm{cPRA}}}
\newcommand{\hicpra}{\ensuremath{\mathrm{hi\text{-}cPRA}}}
\newcommand{\kdpi}{\ensuremath{\mathrm{KDPI}}}
\newcommand{\cit}{\ensuremath{\mathrm{CIT}}}
\newcommand{\ilp}{\ensuremath{\mathrm{ILP}}}
\newcommand{\moea}{\ensuremath{\mathrm{MOEA}}}
\newcommand{\ucbone}{\ensuremath{\mathrm{UCB1}}}
\newcommand{\linucb}{\ensuremath{\mathrm{LinUCB}}}
\newcommand{\frrmab}{\ensuremath{\mathrm{FRRMAB}}}
\newcommand{\picef}{\ensuremath{\mathrm{PICEF}}}
\newcommand{\anzkx}{\ensuremath{\mathrm{ANZKX}}}

\newcommand{\resultfig}[2]{%
  \IfFileExists{#1}{\includegraphics[width=#2]{#1}}%
  {\fbox{\parbox[c][3cm][c]{#2}{\centering\small\tbd{figure: #1}}}}}

\hyphenation{Ad-e-laide Auck-land Bris-bane}

%% ============================================================
\begin{document}

\title{Dynamic Kidney Exchange Optimization Under Realistic
Clinical Constraints for Equitable Allocation}

\author{%
\thanks{Neha and Y.~Amer are with the School of Computer and
Mathematical Sciences, University of Adelaide, Adelaide, SA 5005,
Australia (e-mail: nehar5383@gmail.com). }}

\maketitle
\begin{abstract}
Kidney exchange programs clear their pools of incompatible
donor--patient pairs every month, but fairness in kidney exchange has
mostly been studied one clearing round at a time. We examine how the
price of fairness evolves over repeated rounds, whether myopic or
batched clearing better preserves the balance between efficiency and
equity, and what a multi-objective search adds to an exact
optimization baseline. We use a synthetic simulator informed by public
aggregate statistics of the Australian and New Zealand Paired Kidney
Exchange, with Poisson arrivals, withdrawal after 6, 12 or 18 months,
exits to deceased-donor transplantation and cold ischemia limits
across six regional nodes. Each round is a four-objective problem
covering transplant volume, waiting-time inequality, access for highly
sensitized patients and logistics disparity, solved by an ILP-guided
NSGA-III whose mutation rate is chosen by a UCB1 or LinUCB bandit.
Over 25 rounds and 100 paired runs, myopic clearing with this
framework transplanted as many patients as the volume-maximizing ILP,
while reducing the access gap for highly sensitized patients by
21--24\%, the logistics gap by 28--34\% and waiting-time inequality by
3--4\%. The price of fairness averaged 1.6--1.9\% per round, stayed
flat after the first round, and 94--100\% of the deferred transplants
were made in later rounds. Batched clearing never cleared at the
six-month threshold and lost 12--24\% of transplants at the longer
ones. An exact two-stage ILP, which selects the fairest matching among
those of maximum volume, was more effective at maximum volume: it made
1.0--1.3 more transplants per run with equal or lower inequality. The
evolutionary front is useful below maximum volume: an 80\% floor
lowered the access gap by 17--23\% relative to the two-stage ILP for
2.1--2.9\% fewer transplants. All results are from simulation, not
patient-level clinical data.
\end{abstract}

\begin{IEEEkeywords}
Kidney exchange, dynamic matching, price of fairness,
multi-objective evolutionary optimization, NSGA-III, integer
programming, adaptive operator selection, highly sensitized patients.
\end{IEEEkeywords}

%% ============================================================
%% SECTION 1 — INTRODUCTION
%% ============================================================
\section{Introduction}
\label{sec:intro}

\IEEEPARstart{K}{idney} exchange programs (KEPs) enable living-donor
transplants between biologically incompatible donor--patient pairs by
arranging cyclic exchanges and chains initiated by altruistic
donors~\cite{roth2004kidney}.
In August 2019, New Zealand joined the Australian Paired Kidney Exchange
(AKX) to form the trans-Tasman Australian and New Zealand Paired Kidney
Exchange (ANZKX), which involves more than 25 transplant units and
matches pairs continuously through the national OrganMatch system as
new pairs enter the program~\cite{anzkxprotocol2021}.
Between 60 and 90 pairs are registered at any one time, and the
program performs 50--66 transplants per year~\cite{ANZKX2021,DonateLife2024}.

Operational KEPs rely on combinatorial optimization.
Selecting a maximum set of disjoint exchanges is NP-hard when cycles of
length $K \geq 3$ are allowed~\cite{abraham2007clearing}, and integer
linear programs (ILPs), such as the position-indexed chain-edge
formulation (PICEF)~\cite{dickerson2016picef}, are the standard way of
clearing a pool~\cite{biro2021modelling}.
These methods optimize the number of transplants, or a hierarchy of
criteria usually led by it, over one static compatibility graph.

Fairness has mostly been studied in the same static setting.
Dickerson~\textit{et~al.}~\cite{dickerson2014price} defined the price of
fairness (PoF) in kidney exchange as the relative loss of transplants
caused by favouring highly sensitized patients, whose high calculated
panel-reactive antibody (cPRA) levels leave them with few compatible
donors.
Lexicographic and individual notions of fairness
followed~\cite{mcelfresh2018balancing,farnadi2021individual}.
These studies define fairness among patients, and its cost, within a
clearing round.
Real programs, however, clear the pool repeatedly: unmatched patients
keep waiting, some withdraw, some receive a deceased-donor kidney, and
new pairs arrive.

The dynamic side of kidney exchange has been studied mainly in terms of
efficiency.
\"{U}nver~\cite{unver2010dynamic} modeled kidney exchange as a
stochastic dynamic process.
Akbarpour~\textit{et~al.}~\cite{akbarpour2020thickness} showed that the
value of waiting for a thicker pool depends on information about
departures: when departures can be anticipated, waiting until pairs
become critical outperforms greedy matching, whereas without such
information greedy matching is close to optimal.
FutureMatch~\cite{dickerson2015futurematch} combines human value
judgements with machine learning to match in dynamic environments.
These models evaluate dynamic clearing in terms of efficiency; in this
paper, we track fairness and the efficiency given up for it across
sequential clearing rounds.

\begin{figure}[!t]
  \centering
  \includegraphics[width=\columnwidth]{fig1_anzkx.jpg}
  \caption{Dynamic pool model.
           (a)~Six regional nodes used to represent ANZKX transplant
           units and inter-center kidney transport (routes shown are
           illustrative; under the baseline 24-hour cold ischemia
           limit assumed in this study, all routes are feasible).
           (b)~Monthly timeline: pairs arrive as a Poisson process
           ($\lambda = 10$ per month) and each patient leaves for a
           deceased-donor transplant with probability $\mu_d = 0.02$
           per month; a pair withdraws after waiting
           $\tau \in \{6, 12, 18\}$ months, counted from its arrival.}
  \label{fig:anzkx}
\end{figure}

In earlier work, we developed a fairness-aware multi-objective
evolutionary framework for ANZKX with three objectives and measured a
PoF of 9.6\% under static pool conditions.
Three limitations motivate the present study.
First, the pool was static: pairs neither arrived nor left between
clearing rounds.
Second, the competing exit to the deceased-donor waiting list, which
reduces patients from the living-donor pool, was not modeled.
Third, the cold ischemia time (CIT) limit on transporting a kidney
between geographically dispersed transplant centers was ignored.
In ANZKX, kidneys are shipped by courier between donor and recipient
hospitals, and the CIT usually exceeds 10 hours for trans-Tasman
exchanges and for exchanges between the east and west coasts of
Australia~\cite{anzkxprotocol6}.

This paper addresses all three limitations (Fig.~\ref{fig:anzkx}).
The central research question is: \emph{how does the price of fairness
evolve across sequential dynamic clearing rounds in an ANZKX-calibrated
simulator under realistic operational constraints, and which clearing
strategy, myopic or batched, best preserves the efficiency-equity
trade-off over time?}
We address six research questions: how $\pof$ evolves over rounds
(RQ1); how myopic and batched clearing compare (RQ2); how the
framework compares with the baselines, including an exact two-stage
ILP that picks the fairest maximum matching (RQ3); what the
evolutionary front offers below maximum volume (RQ4); what operator
selection and ILP guidance contribute (RQ5); and how sensitive and
scalable the results are (RQ6).

The main contributions of this paper are as follows.
\begin{enumerate}
  \item \textbf{A price-of-fairness trajectory.}
    We track the price of fairness over 25 dynamic clearing rounds
    instead of a single round. The fairness-aware framework gives up
    about 2\% of transplants per round on average, yet over the whole
    period it transplants as many patients as the volume-maximizing
    ILP, because most deferred transplants are made in later rounds.
  \item \textbf{Clearing strategies compared on fairness.}
    We compare myopic and batched clearing on all four objectives, not
    only on transplant volume. With withdrawals and competing exits,
    myopic clearing transplants more patients, including more highly
    sensitized patients, at every withdrawal threshold studied.
  \item \textbf{An ANZKX-calibrated dynamic simulator} with Poisson
    arrivals, virtual crossmatching, withdrawal after $\tau$ months, a
    deceased-donor competing exit, and CIT limits between six regional
    nodes, using common random numbers so that every method faces the
    same patients.
  \item \textbf{An ILP-guided NSGA-III, assessed against an exact
    benchmark.} Seeding the search with the volume-optimal ILP
    matching, giving its structures priority in the repair, and
    keeping it in an external archive reduces search-failure bias in
    the price of fairness. We compare the framework with an exact
    two-stage ILP at three levels (identical pools, fixed
    trajectories and full trajectories). The exact method is more
    effective at maximum volume; the evolutionary front provides
    trade-off solutions when some volume is given up.
  \item \textbf{A broad evaluation} of ten configurations over
    three withdrawal thresholds with 100 paired runs each, including
    sensitivity, scalability, selection-rule and out-of-distribution
    tests on PrefLib dataset~00036~\cite{mattei2013preflib}.
\end{enumerate}
Section~\ref{sec:related} reviews related work.
Section~\ref{sec:problem} presents the dynamic pool model, the four
objectives and the ILP baselines.
Section~\ref{sec:method} describes the framework.
Section~\ref{sec:experiments} details the experimental setup, and
Section~\ref{sec:results} reports the results.
Section~\ref{sec:conclusion} concludes.

%% ============================================================
%% SECTION 2 — RELATED WORK  (unchanged)
%% ============================================================
\section{Related Work}
\label{sec:related}

Kidney exchange research draws on combinatorial optimization, market
design, and evolutionary computation.
We review four streams of work, and Table~\ref{tab:related} summarizes
how representative studies cover the design criteria of our framework.

\subsection{Integer Programming for Kidney Exchange}
\label{sec:related:ilp}

Abraham~\textit{et~al.}~\cite{abraham2007clearing} proved that clearing
a kidney exchange pool is NP-hard when cycles of length $K \geq 3$ are
allowed, and proposed a branch-and-price algorithm that made nationwide
clearing practical.
Roth~\textit{et~al.}~\cite{roth2007efficient} had shown that three-way
exchanges substantially increase the number of transplants, while
longer exchanges add little.
Saidman~\textit{et~al.}~\cite{saidman2006increasing} reached a similar
conclusion on clinical data, and
Montgomery~\textit{et~al.}~\cite{montgomery2006domino} introduced domino
paired donation, in which a non-directed donor starts a chain.

Formulations have since become stronger and more compact.
Constantino~\textit{et~al.}~\cite{constantino2013new} proposed the first
compact formulations, whose size grows polynomially with the instance.
Glorie~\textit{et~al.}~\cite{glorie2014kidney} noted that pools with many
hard-to-match patients often need long cycles and chains, and gave an
iterative branch-and-price algorithm for multiple criteria with a
polynomial-time pricing problem.
Anderson~\textit{et~al.}~\cite{anderson2015new} adapted travelling
salesman techniques to find long chains in large instances, and
Dickerson~\textit{et~al.}~\cite{dickerson2016picef} introduced
position-indexed formulations, including PICEF for pools with both
cycles and chains.
More recently, Riascos-\'{A}lvarez~\textit{et~al.}~\cite{riascos2023branch}
enhanced branch-and-price with decision diagrams, and Lam and
Mak-Hau~\cite{lam2020branch} developed a branch-and-cut-and-price
algorithm for the cardinality-constrained multi-cycle problem.
Two surveys compare these models~\cite{barkel2026operational,biro2021modelling}.
In practice, ANZKX matches pairs through the national OrganMatch
system~\cite{anzkxprotocol2021}.

These methods optimize the number of transplants, or a hierarchy of
criteria usually led by it~\cite{biro2021modelling}, within a single
round.
We instead treat fairness and efficiency as separate objectives and
trade them off with a multi-objective evolutionary approach.

\subsection{Dynamic Kidney Exchange}
\label{sec:related:dynamic}

Pairs enter an exchange over time, and ANZKX now matches them
continuously as they register~\cite{ANZKX2021}.
\"{U}nver~\cite{unver2010dynamic} modeled kidney exchange as a dynamic
process and characterized efficient mechanisms for two-way and
multi-way exchanges.
Dickerson~\textit{et~al.}~\cite{dickerson2012optimizing} showed that long
chains help little in large random pools but help in small, highly
sensitized ones; with monthly matching, a chain cap of four performed
best, which motivates our choice of $L = 4$.
FutureMatch~\cite{dickerson2015futurematch} combines human value
judgments with machine learning to learn edge weights and vertex
potentials for dynamic matching.
Glorie~\textit{et~al.}~\cite{glorie2022health} embedded a Markov model
of patient health in a discrete-event simulation of the Dutch program
and compared allocation policies by quality-adjusted life-years; a
policy that maximized health value captured more of the potential gain
than the current policy, and systematically enrolling altruistic
donors reduced inequity for patients with blood type O.

Two results bear directly on when to clear a pool.
Anderson~\textit{et~al.}~\cite{anderson2017efficient} showed that greedy
matching is near optimal for average waiting time among a large class
of policies that includes batching, and that three-way cycles
substantially reduce waiting time.
Akbarpour~\textit{et~al.}~\cite{akbarpour2020thickness} showed that
waiting for a thicker pool beats greedy matching only when departures
can be anticipated.
Both results concern efficiency alone; whether they hold once fairness
enters the clearing decision is part of our question.
In a more general setting, Huang and Shu~\cite{huang2021online} analyzed
online stochastic matching with Poisson arrivals through a natural
linear program.

Exchanges across hospitals or countries raise further problems.
Ashlagi and Roth~\cite{ashlagi2014free} showed that hospitals may enroll
only hard-to-match pairs, yet full participation is individually
rational in almost every large pool.
Hajaj~\textit{et~al.}~\cite{hajaj2015strategy} proposed a strategy-proof
and efficient credit mechanism with significant gains on US data.
For international programs,
Klimentova~\textit{et~al.}~\cite{klimentova2021fairness} divided
transplants fairly among agents over time,
Benedek~\textit{et~al.}~\cite{benedek2025partitioned} modeled each
round as a partitioned matching game, and
Druzsin~\textit{et~al.}~\cite{druzsin2024enckep} used the ENCKEP
simulator to evaluate programs over five years.
We instead optimize several fairness criteria for patients across
consecutive rounds.

\subsection{Fairness in Kidney Exchange}
\label{sec:related:fairness}

A patient's chance of a match depends strongly on blood group and
sensitization, and highly sensitized patients are among the hardest to
match~\cite{ANZKX2021}.
Bertsimas~\textit{et~al.}~\cite{bertsimas2011price} defined the price of
fairness as the relative loss of efficiency when a fair allocation, such
as a max-min or proportionally fair one, replaces an efficient one, and
Caragiannis~\textit{et~al.}~\cite{caragiannis2012efficiency} quantified
this loss for proportional, envy-free, and equitable division.

Dickerson~\textit{et~al.}~\cite{dickerson2014price} brought the price of
fairness to kidney exchange, defining highly sensitized patients by
$\text{cPRA} \geq 80$.
They proved it small in the standard theoretical model but found it
often far from zero on real data, reaching 19\% in some runs.
McElfresh and Dickerson~\cite{mcelfresh2018balancing} prioritized highly
sensitized patients while bounding the loss of efficiency.
Freedman~\textit{et~al.}~\cite{freedman2020adapting} turned elicited
human judgements into patient weights, and
Farnadi~\textit{et~al.}~\cite{farnadi2021individual} used multiple
optimal solutions to give patients equal chances of a transplant.
St-Arnaud~\textit{et~al.}~\cite{starnaud2025adaptation} and
Chang~\textit{et~al.}~\cite{chang2025optimizing} built lotteries over
exchange plans, the latter with polynomial-time guarantees for
objectives such as Nash social welfare.
Zhang~\textit{et~al.}~\cite{zhang2026fairness} required outcomes to be
independent of protected features given sensitization.
At program level, Bir\'{o}~\textit{et~al.}~\cite{biro2021modelling}
show that European programs use different optimization criteria.

Multi-agent programs track fairness among countries across
rounds~\cite{klimentova2021fairness}, but the patient-level studies
above measure it within a single round.
We treat waiting-time equity, access for highly sensitized patients,
and logistics disparity as separate objectives and follow them across
rounds.

\subsection{Multi-Objective Evolutionary Algorithms and Adaptive
Operator Selection}
\label{sec:related:moea}

NSGA-II~\cite{deb2002nsga2} set the standard for Pareto-based search.
NSGA-III~\cite{deb2014nsga3} extended it to four or more objectives
through reference points, which keep the population spread across the
front.
MOEA/D~\cite{zhang2007moead} is the main decomposition-based
alternative, and Trivedi~\textit{et~al.}~\cite{trivedi2017survey} survey
its extensions.
Ishibuchi~\textit{et~al.}~\cite{ishibuchi2017reference} showed how the
reference point affects hypervolume comparisons and how to choose it
fairly.
\begin{table*}[!htbp]
\caption{Comparison of representative kidney exchange studies against
  key design criteria. \checkmark~= criterion explicitly addressed;
  {--}~= absent or not explicit.
  Groups: (A)~static ILP, (B)~dynamic KEP, (C)~fairness-aware KEP.}
\label{tab:related}
\centering
\small
\setlength{\tabcolsep}{5pt}
\begin{tabular}{@{}llcclcc@{}}
\toprule
\textbf{Work} & \textbf{Year} & \textbf{Dyn.\ Pool} &
  \textbf{Obj.} & \textbf{Fairness Dimension} &
  \textbf{AOS} & \textbf{hi-cPRA} \\
\midrule
\multicolumn{7}{@{}l}{\textit{(A) Static ILP Formulations}} \\[2pt]
Abraham~\textit{et~al.}~\cite{abraham2007clearing}
  & 2007 & --          & 1 & None                        & -- & -- \\
Constantino~\textit{et~al.}~\cite{constantino2013new}
  & 2013 & --          & 1 & None                        & -- & -- \\
Anderson~\textit{et~al.}~\cite{anderson2015new}
  & 2015 & --          & 1 & None                        & -- & -- \\
Dickerson~\textit{et~al.}~\cite{dickerson2016picef}
  & 2016 & --          & 1 & None                        & -- & -- \\
Riascos-\'{A}lvarez~\textit{et~al.}~\cite{riascos2023branch}
  & 2024 & --          & 1 & None                        & -- & -- \\
\midrule
\multicolumn{7}{@{}l}{\textit{(B) Dynamic KEP}} \\[2pt]
\"{U}nver~\cite{unver2010dynamic}
  & 2010 & \checkmark  & 1 & None                        & -- & -- \\
Dickerson~\textit{et~al.}~\cite{dickerson2012optimizing}
  & 2012 & \checkmark  & 1 & None                        & -- & -- \\
Dickerson and Sandholm~\cite{dickerson2015futurematch}
  & 2015 & \checkmark  & 1 & Human value judgements      & -- & -- \\
Anderson~\textit{et~al.}~\cite{anderson2017efficient}
  & 2017 & \checkmark  & 1 & None                        & -- & -- \\
Akbarpour~\textit{et~al.}~\cite{akbarpour2020thickness}
  & 2020 & \checkmark  & 1 & None                        & -- & -- \\
\midrule
\multicolumn{7}{@{}l}{\textit{(C) Fairness-Aware KEP}} \\[2pt]
Dickerson~\textit{et~al.}~\cite{dickerson2014price}
  & 2014 & --          & 1 & Lexicographic, weighted     & -- & \checkmark \\
Freedman~\textit{et~al.}~\cite{freedman2020adapting}
  & 2018 & --          & 1 & Elicited human values       & -- & -- \\
McElfresh and Dickerson~\cite{mcelfresh2018balancing}
  & 2018 & --          & 2 & Hybrid lexicographic        & -- & \checkmark \\
Farnadi~\textit{et~al.}~\cite{farnadi2021individual}
  & 2021 & --          & 1 & Individual                  & -- & -- \\
Zhang~\textit{et~al.}~\cite{zhang2026fairness}
  & 2025 & --          & 1 & Calibration-based           & -- & -- \\
\midrule
\textbf{This work}
  & \textbf{2026} & \checkmark & \textbf{4}
  & \textbf{Waiting-time Gini, hi-cPRA gap, logistics}
  & \checkmark & \checkmark \\
\bottomrule
\end{tabular}
\end{table*}
Adaptive operator selection (AOS) chooses among operators, or values of
an operator parameter, during the search.
Thierens~\cite{thierens2005adaptive} proposed adaptive pursuit as an
alternative to probability matching.
Bandit methods followed, building on the UCB1 analysis of
Auer~\textit{et~al.}~\cite{auer2002finite}:
Maturana~\textit{et~al.}~\cite{maturana2009extreme} combined
extreme-value credit with dynamic bandits,
Fialho~\textit{et~al.}~\cite{fialho2010analyzing} compared UCB, dynamic,
and sliding-window bandits, and
Li~\textit{et~al.}~\cite{li2014frrmab} proposed FRRMAB for MOEA/D.
Pei~\textit{et~al.}~\cite{pei2025aos} classify AOS methods as stateless
or state-based, and deep Q-networks now drive state-based selection in
multi-objective~\cite{yin2024dueling} and constrained
multi-objective~\cite{ming2024constrained} algorithms.
Contextual bandits such as LinUCB~\cite{li2010contextual} sit between the
two: they use a context vector but need no neural network.

Deep reinforcement learning has also been applied to multi-objective
optimization directly~\cite{li2021deeprl}, but it needs training before
use; bandits adapt online, which suits a pool that changes every month.
Yu~\textit{et~al.}~\cite{yu2025fairness} connect multi-objective
optimization with fairness-aware machine learning, a framing close to
ours.

In this paper, we guide NSGA-III with the volume-optimal ILP matching,
combine it with bandit-based mutation-rate selection, compare
stateless and contextual bandits in fairness-aware dynamic kidney
exchange, and assess the result against an exact two-stage ILP.

%% ============================================================
\section{Problem Formulation}
\label{sec:problem}
%% ============================================================

\subsection{Dynamic Pool Model}
\label{sec:problem:model}

We model the exchange in monthly rounds $t = 1,\dots,T$. The pool
$V_t$ holds incompatible donor--patient pairs and altruistic donors,
and $P_t \subseteq V_t$ denotes the patients. Each patient $i$ has a
blood type, a cPRA value, a waiting time $w_i$ in months and a home
center $h(i)$: Perth, Adelaide, Melbourne, Sydney, Brisbane or
Auckland. Highly sensitized patients form the set
$S_t = \{ i \in P_t : \cpra_i \geq 0.80 \}$.

Once the matching $M_t$ of round $t$ is implemented, the pool evolves
as
\begin{equation}
  V_{t+1} = \bigl(V_t \setminus M_t \setminus X_t \setminus W_t\bigr)
            \cup A_t ,
  \label{eq:pool}
\end{equation}
in this order. Matched pairs leave. Each remaining patient leaves for
a deceased-donor transplant with probability $\mu_d$ (the competing
exit $X_t$), since patients registered in the Australian paired
donation program are often also listed for a deceased-donor
kidney~\cite{Ferrari2012}. Waiting times then grow by one month, pairs that have
waited $\tau$ months withdraw ($W_t$), and
$|A_t| \sim \mathrm{Poisson}(\lambda)$ new pairs arrive. The
withdrawal threshold $\tau$ is a modeling assumption informed
by~\cite{Ferrari2012}, and we vary it over $\{6, 12, 18\}$ months.

New pairs follow the Saidman model~\cite{saidman2006increasing} with
three changes. Blood types follow Australian population frequencies.
Thirty percent of patients have $\cpra \in [0.80, 0.99]$, and the rest
are split equally between $[0, 0.05)$ and $[0.05, 0.80)$. A pair
joins the pool only if its donor cannot give to its own patient,
because of ABO incompatibility or a positive crossmatch. Eight percent
of arrivals are altruistic donors.

\subsection{Compatibility and Logistics}
\label{sec:problem:compat}

The compatibility graph $G_t = (V_t, E_t)$ has an arc $(i, j)$ when
the donor of $i$ can give to the patient of $j$:
\begin{equation}
\begin{aligned}
  (i,j) \in E_t \iff{} & \mathrm{ABO}(i,j) \wedge \mathrm{XM}^{-}(i,j) \\
  & \wedge\; T^{\mathrm{tr}}_{h(i)h(j)} + T^{\mathrm{prep}} \leq \cit_{\max} .
\end{aligned}
\label{eq:edge}
\end{equation}
The virtual crossmatch $\mathrm{XM}^{-}(i,j)$ is negative with
probability $1 - \cpra_j$. It is drawn once for each donor--patient
combination, so the same couple gets the same result in every round
and configuration. We model the transport time $T^{\mathrm{tr}}_{ab}$
between two different centers as the scheduled flight time, plus one
hour of ground transport at each end, plus a wait for a suitable
flight (2~h domestic, 4~h trans-Tasman); within a center it is zero.
We assume $T^{\mathrm{prep}} = 6$~h and $\cit_{\max} = 24$~h. These
transport and CIT values are modeling assumptions, not ANZKX
operational parameters. HLA profiles are drawn at random in this
simplified model, and each arc carries an HLA weight
$\omega_{ij} = 1 - m_{ij}/6$, where $m_{ij}$ is the number of HLA
mismatches; it is used only to break ties.
\begin{figure}[!t]
\centering
\includegraphics[width=\columnwidth]{fig2_structures.jpg}
\caption{Exchange structures. Each box is a donor (D) and patient
(P) pair, and an arrow means the donor of one pair gives a kidney to
the patient of the next. (a)~Two-way and (b)~three-way cycles; we
allow cycles of up to $K = 3$ pairs. (c)~A chain started by an
altruistic donor (A) with $L = 4$ transplants; the donor of the last
pair (grey) does not donate, and the chain ends there.}
\label{fig:structures}
\end{figure}
A matching consists of vertex-disjoint cycles of at most $K = 3$
pairs and chains of at most $L = 4$ transplants, each started by an
altruistic donor. Fig.~\ref{fig:structures} illustrates the three
types of exchange. We enumerate all cycles and chains once per round,
and every method works on the same list. To keep the enumeration
tractable, we keep at most 5{,}000 chains per round: all chains with
one or two transplants and a random sample of longer ones. When this
cap binds, we add the chains of the exact optimum found with
PICEF~\cite{dickerson2016picef}, so the list always contains a
volume-optimal matching.

\subsection{Objective Functions}
\label{sec:problem:objectives}

For a matching $M$, let $P = P_t$ and $S = S_t$, and let $R$ and $U$
be the patients at the distant centers (Perth, Auckland) and the
central ones (Adelaide, Melbourne, Sydney, Brisbane). All four sets
contain patients only; altruistic donors are not counted.

\begin{definition}[Transplant volume]
$F_1(M) = |M \cap P|$, the number of patients transplanted.
\end{definition}

\begin{definition}[Waiting-time equity]
$F_2(M)$ is the Gini coefficient of the waiting times after the round,
where matched patients keep their waiting time $w_i$ and the others move
on to $w'_i = w_i + 1$:
\begin{equation}
  F_2(M) = \frac{\sum_{i \in P}\sum_{j \in P} |w'_i - w'_j|}
                {2 |P|^2 \, \overline{w'}} ,
  \label{eq:gini}
\end{equation}
where $\overline{w'}$ is the mean of $w'$, and $F_2 = 0$ when
$|P| < 2$ or all waiting times are zero.
\end{definition}

\begin{definition}[High-cPRA access gap]
\begin{equation}
  F_3(M) = \left| \frac{|M \cap S|}{|S|} - \frac{|M \cap P|}{|P|}
  \right| ,
  \label{eq:f3}
\end{equation}
with $F_3 = 0$ when $S = \emptyset$.
\end{definition}
$F_3$ compares the match rate of highly sensitized patients with the
overall match rate; the absolute value penalizes both neglect and
over-correction. A smaller $F_3$ therefore means more equal match
rates, not necessarily more transplants for highly sensitized
patients.

\begin{definition}[Logistics disparity]
\begin{equation}
  F_4(M) = \left| \frac{|M \cap R|}{|R|} - \frac{|M \cap U|}{|U|}
  \right| ,
  \label{eq:f4}
\end{equation}
with $F_4 = 0$ when $R$ or $U$ is empty.
\end{definition}
$F_4$ is a modeled proxy for geographic disparity, not a validated
measure of inequity.

Each round then solves
\begin{equation}
  \min_{M \in \mathcal{M}_t} \;
  \bigl( 1 - F_1/|P|,\; F_2,\; F_3,\; F_4 \bigr) ,
  \label{eq:mop}
\end{equation}
where $\mathcal{M}_t$ is the set of feasible matchings, and all four
objectives lie in $[0,1]$. The empty matching has $F_3 = F_4 = 0$ and
is Pareto-optimal, so it is never implemented.

\subsection{Price of Fairness}
\label{sec:problem:pof}

\begin{definition}[Price of fairness]
Let $F_1^{\ilp}(t)$ be the largest number of transplants possible in
round $t$ and $F_1(t)$ the number in the implemented matching. Then
\begin{equation}
  \pof = \frac{F_1^{\ilp}(t) - F_1(t)}{F_1^{\ilp}(t)} \times 100\% ,
  \label{eq:pof}
\end{equation}
which is undefined when $F_1^{\ilp}(t) = 0$.
\end{definition}
We summarize the trajectory $\{\pof\}_{t=1}^{T}$ by its mean over the
rounds where it is defined, its variance and its per-run slope.

$\pof$ compares two matchings of the same pool. Over a trajectory,
however, policies produce different pools, and a patient passed over
in one round may be matched in the next. We therefore also measure
the loss over the whole trajectory,
\begin{equation}
  \Delta = 1 - \frac{\sum_t F_1(t)}{\sum_t F_1^{\mathrm{MI}}(t)} ,
  \label{eq:delta}
\end{equation}
where $F_1^{\mathrm{MI}}(t)$ is the number of transplants made by the
myopic ILP in the run with the same seed. Unlike $\pof$, $\Delta$ can
also compare clearing policies, since batched clearing skips rounds.

\subsection{Integer Programming Baselines}
\label{sec:problem:ilp}

Let $x_k$ be a binary variable for each enumerated structure $k$ with
$|k|$ transplants. The volume-maximizing ILP is
\begin{equation}
  \max \sum_k \bigl(|k| + \epsilon\, \bar\omega_k\bigr) x_k
  \quad \text{s.t.} \quad
  \sum_{k \ni v} x_k \leq 1 \;\; \forall v \in V_t ,
  \label{eq:ilp}
\end{equation}
where $\bar\omega_k$ is the normalized HLA weight of structure $k$ and
$\epsilon = 10^{-3}$ breaks ties without giving up a transplant. The
lexicographic baseline of
Dickerson~\textit{et~al.}~\cite{dickerson2014price} first maximizes
the number of highly sensitized patients transplanted, then the total
number of transplants.

\emph{Two-stage ILP.} Stage~1 is~\eqref{eq:ilp} and gives
$F_1^{*} = F_1^{\ilp}(t)$ with its matching $M^{\ilp}_t$. Stage~2 keeps
the same structures and packing constraints, fixes
$\sum_k |k|\, x_k = V$ with $V = F_1^{*}$, and minimizes
\begin{equation}
  \frac{F_2}{r_2} + \frac{F_3}{r_3} + \frac{F_4}{r_4}
  - \epsilon' \sum_k \bar\omega_k x_k ,
  \qquad r_k = \max\bigl(F_k(M^{\ilp}_t),\, 0.01\bigr) ,
  \label{eq:twostage}
\end{equation}
so each fairness objective is scaled by its value in the
volume-maximizing matching (the floor avoids division by values near
zero), and $\epsilon' = 10^{-7}$ breaks ties by HLA quality. With $V$
fixed, all three objectives are linear. $F_3$ and $F_4$ are absolute
differences of match rates whose overall rate $V/|P|$ is a constant,
so one auxiliary variable each linearizes them. For $F_2$, the sum
$\sum_i w'_i = \sum_i w_i + |P| - V$ is constant, so the denominator
of~\eqref{eq:gini} is fixed. For two patients with $w_i \neq w_j$, the
sign of $w'_i - w'_j$ does not depend on the matching, because
$|w_i - w_j| \geq 1$ and the matching changes each waiting time by at
most one; these terms are linear. Pairs with $w_i = w_j$ need one
auxiliary variable each. The formulation is therefore exact. We call
this baseline the \emph{two-stage ILP (sum)}. A second variant,
\emph{two-stage ILP ($F_3$)}, minimizes $F_3$ alone, with
$10^{-5}\, F_4/r_4$ and the HLA term as tie-breakers. To measure the
fairness available below maximum volume (Section~\ref{sec:exp:levels}),
we also solve stage~2 with $V < F_1^{*}$.

All ILPs are solved with CBC, and HiGHS is used if CBC fails. Stage~2
has a time limit of 60~s per round in the trajectories; when it is
reached, the best matching found is used, and we report how often this
happened.

%% ============================================================
\section{Proposed Framework}
\label{sec:method}
%% ============================================================

\subsection{Overview}
\label{sec:method:overview}

Fig.~\ref{fig:framework} shows one round of the framework, and
Algorithm~\ref{alg:main} gives the full trajectory. The framework
couples a clearing policy, which decides \emph{when} the pool is
cleared, with an ILP-guided NSGA-III, which decides \emph{how} it is
cleared. In each round in which the policy clears, the
ILP~\eqref{eq:ilp} first gives the maximum number of transplants
$F_1^{\mathrm{ILP}}(t)$. NSGA-III then approximates the Pareto front
of~\eqref{eq:mop}, one matching from the front is implemented, and the
pool evolves by~\eqref{eq:pool}. Within a round, a bandit sets the
mutation rate at every generation, so the policy acts once per round
and the bandit once per generation.

\begin{figure*}[!t]
\centering
\includegraphics[width=\textwidth]{fig3_framework.jpg}
\caption{One round of the proposed framework. The ILP solution seeds
NSGA-III, whose mutation rate is chosen by a bandit; rounds with at
most 200 feasible matchings are solved exactly. The implemented
matching keeps at least 90\% of the maximum transplants and is
closest to the ideal point in $(F_2, F_3, F_4)$.}
\label{fig:framework}
\end{figure*}

\begin{algorithm}[!htbp]
\caption{One dynamic trajectory}
\label{alg:main}
\begin{algorithmic}[1]
\Require initial pool $V_1$; policy $\pi \in \{\pi_M, \pi_B\}$;
  bandit $\mathcal{B}$; rounds $T$; population size $N$;
  generations $G_{\max}$
\State $\mathcal{F}_0 \gets \emptyset$ \Comment{previous front}
\For{$t = 1, \dots, T$}
  \State $M_t \gets \emptyset$
  \If{$\pi$ clears in round $t$}
    \State build $G_t$ and the structure list $\mathcal{K}_t$
    \State $(M^{\mathrm{ILP}}_t, F_1^{\mathrm{ILP}}(t)) \gets$
           solve~\eqref{eq:ilp}
    \If{$\mathcal{K}_t$ admits at most $2N$ matchings}
      \State $\mathcal{F}_t \gets$ exact Pareto front by enumeration
    \Else
      \State $\mathcal{Q} \gets \{M^{\mathrm{ILP}}_t\} \cup
             \mathcal{F}_{t-1} \cup$ random individuals
      \For{$g = 1, \dots, G_{\max}$}
        \State $p_m \gets \mathcal{B}.\mathrm{select}(\mathbf{x}_t)$
        \State create offspring by uniform crossover and bit-flip
               mutation with rate $p_m$
        \State repair; apply NSGA-III environmental selection
        \State update the archive; $\mathcal{B}.\mathrm{update}(r_g)$
      \EndFor
      \State $\mathcal{F}_t \gets$ archive
    \EndIf
    \State $M_t \gets \mathrm{select}(\mathcal{F}_t)$; record
           $\mathrm{PoF}_t$
  \EndIf
  \State $V_{t+1} \gets$ update the pool by~\eqref{eq:pool}
\EndFor
\end{algorithmic}
\end{algorithm}

\subsection{Encoding, Repair and Variation}
\label{sec:method:nsga3}

A solution is a binary vector with one bit per vertex, indicating
whether the pair or altruistic donor is eligible to be matched.
A structure in $\mathcal{K}_t$ is usable only if all of its vertices are
eligible.
The repair operator scans the usable structures in a fixed priority
order, placing the structures of the ILP solution first
(Section~\ref{sec:method:hybrid}), then structures with more
transplants, and then structures with higher HLA quality, and accepts
each structure that is disjoint from those already accepted.
The repaired vector is written back into the population (Lamarckian
repair), so every individual encodes a feasible matching over the same
structure list as the ILP.

NSGA-III~\cite{deb2014nsga3} evolves a population of $N = 100$
individuals for $G_{\max} = 150$ generations per round, using
Das--Dennis reference directions with $p = 12$ divisions.
This yields 455 directions for four objectives; since $N < 455$,
niching associates individuals with their nearest directions.
Offspring are produced by uniform crossover with probability
$p_c = 0.9$ and bit-flip mutation with rate $p_m$, which the bandit
selects at every generation.
Duplicate individuals are removed.

\subsection{ILP Guidance, Archive and Exact Small Rounds}
\label{sec:method:hybrid}

Packing cycles and chains is NP-hard~\cite{abraham2007clearing}, and
NSGA-III alone rarely reaches the largest matching: in 30 pools of 75
vertices it found 32.6 transplants on average against an ILP optimum
of 42.3 and never reached the optimum (Table~\ref{tab:hybrid}). Since
$\pof$ is measured against the ILP optimum, such search failure would
inflate it. We therefore guide the search with the ILP solution: it is
inserted into the initial population, its structures take priority in
the repair, and an external archive collects every non-dominated
solution of the whole population at every generation, so niching
cannot discard the volume-optimal matching. The archive forms the
final front and always contains the ILP optimum, so $\pof$ is measured
from the selected fairness-aware solution against the true volume
optimum. This reduces the bias that search failure would otherwise add
to $\pof$, but it does not prove that all of the remaining loss is
caused by fairness alone: part of it is set by the selection rule
(Section~\ref{sec:method:select}). The same guidance is used in all
six evolutionary configurations; the ILP baselines are exact and need
none. We treat the hybrid as a supporting component that makes $\pof$
measurable, not as a separate algorithmic contribution.

In dynamic configurations, the matchings on the previous round's front
are also inserted into the initial population after the vertices that
have left the pool are removed.
The ILP solution and these warm-start solutions fill at most half of
the population.
Each remaining individual is generated at random with its own
eligibility rate drawn uniformly from $[0.2, 1]$, which keeps the
population diverse.

After matches, withdrawals and exits, a round may contain few
structures.
If $\mathcal{K}_t$ admits at most $2N = 200$ matchings, including the
empty one, all of them are enumerated and the exact Pareto front is
used.
The bandit is not called in such rounds, because the population could
not be filled with distinct individuals and the reward would carry no
information; its state carries over to the next round.

\subsection{Adaptive Operator Selection}
\label{sec:method:aos}

The bandit chooses among the mutation rates
$\mathcal{A} = \{0.01, 0.05, 0.10, 0.25\}$.
After generation $g$, the chosen arm receives the reward
\begin{equation}
\begin{aligned}
  r_g &= \frac{\Delta\mathrm{HV}_g}{\max_{g' \leq g} \Delta\mathrm{HV}_{g'}}
        \in [0, 1], \\
  \Delta\mathrm{HV}_g &= \max\bigl(0,\, \mathrm{HV}_g - \mathrm{HV}_{g-1}\bigr),
\end{aligned}
\label{eq:reward}
\end{equation}
with $r_g = 0$ while no gain has been observed in the round.
Here $\mathrm{HV}_g$ is the hypervolume of the non-dominated individuals
in the population, computed on the objective vector of~\eqref{eq:mop}
with reference point $(1.1, 1.1, 1.1, 1.1)$.
The reward thus measures the hypervolume gain of a generation relative
to the largest gain so far in the round.

\subsubsection{Stateless UCB1}
UCB1~\cite{auer2002finite} considers the last $\ell = 15$ choices and
selects
\begin{equation}
  a_g = \arg\max_{a \in \mathcal{A}} \;
        \hat{Q}_a + \kappa \sqrt{\frac{2 \ln n}{n_a}} ,
  \label{eq:ucb1}
\end{equation}
where $\hat{Q}_a$ and $n_a$ are the mean reward and the number of plays
of arm $a$ within the window, $n = \sum_{a} n_a$, and $\kappa = 1$.
An arm absent from the window is played first.

\subsubsection{Contextual LinUCB}
LinUCB~\cite{li2010contextual} additionally observes the pool state
\begin{equation}
  \mathbf{x}_t = \Bigl[\tfrac{|P_t|}{150},\;
  \tfrac{|S_t|}{|P_t|},\;
  \tfrac{|X_{t-1}|}{|P_t|},\;
  \tfrac{t}{T},\;
  \mathrm{Gini}_t,\;
  \tfrac{\mathrm{PoF}_{t-1}}{100}\Bigr]^{\top} \in [0,1]^6 ,
  \label{eq:context}
\end{equation}
whose components are the pool size, the share of highly sensitized
patients, the share of competing exits in the previous round, the
progress through the horizon, the waiting-time Gini coefficient of the
pool, and the previous round's price of fairness.
It selects
\begin{equation}
  a_g = \arg\max_{a \in \mathcal{A}} \;
        \boldsymbol{\theta}_a^{\top} \mathbf{x}_t
        + \alpha \sqrt{\mathbf{x}_t^{\top} \mathbf{D}_a^{-1} \mathbf{x}_t} ,
  \label{eq:linucb}
\end{equation}
with $\alpha = 1$, where $\mathbf{D}_a$ and $\boldsymbol{\theta}_a$ are
updated by ridge regression on the observed rewards.
Because $\mathbf{x}_t$ is fixed within a round, LinUCB adapts across
rounds, whereas UCB1 responds to rewards alone.
In dynamic configurations the bandit retains its state across rounds;
in static configurations a new bandit is created in every round and no
warm start is used.

\subsection{Implemented Solution}
\label{sec:method:select}

The empty matching is removed from the final front, and only solutions
with $F_1 \geq 0.9 \max F_1$ are retained.
Among them, the solution closest to the ideal point in
$(F_2, F_3, F_4)$, after min--max scaling over the retained set, is
implemented.
Since the front contains the ILP optimum, $\max F_1 = F_1^{\mathrm{ILP}}(t)$
and
\begin{equation}
  0 \leq \mathrm{PoF}_t \leq 10\% \quad \text{in every cleared round.}
  \label{eq:pofbound}
\end{equation}
Moreover, $\mathrm{PoF}_t = 0$ whenever $F_1^{\mathrm{ILP}}(t) < 10$,
because giving up a single transplant would then exceed the 10\%
allowance, and the rule selects a volume-optimal matching.
This rule differs from the stage-2 objective~\eqref{eq:twostage} of
the two-stage ILP, so the two can choose different maximum matchings.
As the level of $\mathrm{PoF}_t$ depends on the rule, the experiments
also evaluate a 100\% floor (volume-optimal matchings only), an 80\%
floor, and the compromise solution over all non-empty front solutions.

\subsection{Clearing Policies}
\label{sec:method:clearing}

\emph{Myopic} clearing ($\pi_M$) clears in every round.
\emph{Batched} clearing ($\pi_B$) clears only when at least
$n^{*} = 90$ patients are waiting, the upper end of the ANZKX pool
size~\cite{ANZKX2021}.
When the pool cannot reach $n^{*}$, as at $\tau = 6$, $\pi_B$ never
clears.
Two variants are therefore also evaluated: one that additionally clears
once three rounds have passed since the last clearing, and one with
$n^{*} = 60$, the lower end of the ANZKX pool size.

%% ============================================================
\section{Experimental Setup}
\label{sec:experiments}
%% ============================================================

\subsection{Simulation Parameters}
\label{sec:exp:pool}

Table~\ref{tab:params} lists the parameters. We calibrated the
arrival rate by simulating the myopic ILP. With $\lambda = 10$ it
performed 64.9--70.0 transplants per year (ANZKX: 50--66), and the
pool settled at 31.5, 50.3 and 64.0 patients for $\tau = 6$, 12 and 18
(ANZKX: 60--90). No single rate matches both targets, so we also test
$\lambda = 7$, which matches the volume, and $\lambda = 13$, which
matches the pool size. With $\cit_{\max} = 24$~h every center pair is
feasible; the limit first binds at 16~h (Perth--Auckland) and rules
out 73\% of center pairs at 12~h.
Only the high-cPRA share and the ANZKX volume and pool size used for
calibration come from published sources. The deceased-donor exit
rate, the withdrawal thresholds, the transport and CIT constraints and
the simplified HLA model are modeling assumptions, not validated ANZKX
parameters; we therefore vary the most influential of them in the
sensitivity analysis. All results come from synthetic simulation
informed by public ANZKX information; no patient-level clinical data
were used.

\begin{table}[!htbp]
\centering
\caption{Simulation and algorithm parameters (sensitivity values in
parentheses). Basis: P = published evidence, C = calibrated to
published aggregate statistics, A = modeling assumption or design
choice.}
\label{tab:params}
\footnotesize
\setlength{\tabcolsep}{3pt}
\begin{tabular}{@{}p{3.0cm}p{4.0cm}c@{}}
\toprule
Parameter & Value & Basis \\
\midrule
Rounds $T$ / round length & 25 / one month & A \\
Initial pool & 75 vertices, waits $\sim$ U\{0,\dots,5\} & A \\
Arrival rate $\lambda$ & 10 per month (7, 13) & C~\cite{ANZKX2021,DonateLife2024} \\
Withdrawal threshold $\tau$ & 6, 12, 18 months & A~\cite{Ferrari2012} \\
Competing exit rate $\mu_d$ & 0.02 per month (0.01, 0.04) & A~\cite{DonateLife2024} \\
High-cPRA share & 30\% of patients & P~\cite{ANZDATA2026} \\
Altruistic donors & 8\% of arrivals & A \\
Cycle / chain limit & $K = 3$ / $L = 4$ & A \\
$\cit_{\max}$ / $T^{\mathrm{prep}}$, transport & 24~h (16, 12) / 6~h, flight-based & A \\
HLA compatibility & random profiles, tie-break only & A \\
Batching threshold $n^*$ & 90 patients (60) & A~\cite{ANZKX2021} \\
NSGA-III & $N = 100$, $G_{\max} = 150$, $p = 12$ & A \\
Operators & uniform crossover $p_c = 0.9$, bit-flip & A \\
Mutation arms & $\{0.01, 0.05, 0.10, 0.25\}$ & A \\
UCB1 / LinUCB & $\ell = 15$, $\kappa = 1$ / $\alpha = 1$, $d = 6$ & A \\
Selection floor & 90\% (100\%, 80\%, compromise) & A \\
Two-stage ILP limit & 60~s per round (300~s identical pools) & A \\
Runs & 100 paired (30 for sensitivity) & A \\
\bottomrule
\end{tabular}
\end{table}

\subsection{Baselines}
\label{sec:exp:baselines}

Table~\ref{tab:configs} lists the ten configurations we compare.
Configurations 1 and 2 solve each round independently, as in our
earlier static framework. Configurations 3--6 are the proposed
framework under myopic and batched clearing. Configurations 7 and 8
are the volume-maximizing ILP and the lexicographic rule of
Dickerson~\textit{et~al.}~\cite{dickerson2014price}, and 9 and 10 are
the two two-stage ILPs of Section~\ref{sec:problem:ilp}.
\begin{table}[!htbp]
\centering
\caption{Compared configurations. All MOEA configurations use ILP
guidance and the same structure list as the ILPs.}
\label{tab:configs}
\footnotesize
\setlength{\tabcolsep}{3pt}
\begin{tabular}{@{}clllc@{}}
\toprule
\# & Short name & Method & Mutation rate & Clearing \\
\midrule
1 & S-NSGA-II + UCB1  & Static NSGA-II & UCB1         & Myopic \\
2 & S-NSGA-III fixed  & Static NSGA-III                & $p_m = 0.10$ & Myopic \\
3 & D-UCB1 myopic     & Dynamic NSGA-III               & UCB1         & Myopic \\
4 & D-LinUCB myopic   & Dynamic NSGA-III               & LinUCB       & Myopic \\
5 & D-UCB1 batched    & Dynamic NSGA-III               & UCB1         & Batched \\
6 & D-LinUCB batched  & Dynamic NSGA-III               & LinUCB       & Batched \\
7 & Myopic ILP        & ILP~\eqref{eq:ilp}             & --           & Myopic \\
8 & Dickerson lex.    & Lexicographic ILP~\cite{dickerson2014price} & -- & Myopic \\
9 & Two-stage (sum)   & Two-stage ILP~\eqref{eq:twostage} & --        & Myopic \\
10 & Two-stage ($F_3$) & Two-stage ILP, $F_3$ only      & --           & Myopic \\
\bottomrule
\multicolumn{5}{@{}l}{Static: fresh bandit and no warm start in every round.}\\
\multicolumn{5}{@{}l}{Dynamic: bandit state and warm start carried across rounds.}
\end{tabular}
\end{table}

\subsection{Metrics and Statistical Analysis}
\label{sec:exp:metrics}

For each run we record the mean, variance and slope of $\pof$, the
trajectory loss $\Delta$~\eqref{eq:delta}, total and high-cPRA
transplants, high-cPRA withdrawals, $F_2$--$F_4$, hypervolume and the
arms selected. Run $r$ uses the same seed in every configuration, so
all comparisons are paired. A Friedman test is followed by Wilcoxon
signed-rank tests with Bonferroni correction, and effect sizes are
reported as Cohen's $d_z$. The slope of $\pof$ is tested against zero
over all rounds and over rounds 2--25.

\subsection{Comparison with the Two-Stage ILP}
\label{sec:exp:levels}

We compare the framework with the two-stage ILP at three levels,
because each answers a different question.
\begin{itemize}
  \item \emph{Identical pools.} The 30 initial pools of 75 vertices
    are cleared once by every method. This measures, on the same
    instance, how close the framework's matching at maximum volume is
    to the exact optimum of~\eqref{eq:twostage}. Stage~2 has a limit
    of 300~s here.
  \item \emph{Fixed trajectories.} In every round of the framework's
    own trajectories (UCB1, myopic, 30 runs), the two-stage ILP is
    also solved on the same pool, without changing the trajectory.
    This compares the methods on the realistic pools that arise over
    time, round by round.
  \item \emph{Full trajectories.} Each method clears its own pool
    sequence over 25 rounds (100 paired runs). This measures the
    combined effect over time, including how each method shapes
    future pools.
\end{itemize}
To ask what the front offers below maximum volume, we use volume
levels $\rho \in \{1.0, 0.95, 0.90, 0.85, 0.80\}$. Level $\rho$ allows
$V = \lceil \rho F_1^{*} \rceil$ or more transplants. For each level we
record the number of distinct front solutions with $F_1 \geq V$, the
fairest of them, the solution the floor rule would implement, and an
exact reference: stage~2 of the two-stage ILP solved at each level's
volume, taking the best among the volumes from $V$ to $F_1^{*}$. These
level solves use a limit of 120~s on identical pools and 60~s on
trajectory rounds (all three $\tau$, 30 runs). In rounds solved by
enumeration the whole front is exact, so we also count how often
several maximum matchings share the optimal stage-2 score, which
shows how often the deterministic tie-break decides the selected
matching. Finally, the 100\%, 90\% and 80\% floors and the compromise
rule are run over full trajectories at all three $\tau$ (30 runs).

\subsection{Sensitivity, Scalability and PrefLib}
\label{sec:exp:ood}

Sensitivity analyses use 30 runs at $\tau = 12$ and vary the
parameters in parentheses in Table~\ref{tab:params}; the threshold
$n^* = 60$ is also tested at $\tau = 18$. They also test an exit rate
reduced to $0.3\mu_d$ for highly sensitized patients and batched
clearing with a three-round maximum interval, the latter at all
$\tau$. For scalability, we run 10 rounds and 10 runs with initial
pools of 200 vertices (framework and both ILPs) and 500 vertices
(ILPs only). An AOS ablation on 30 pools of 75 vertices compares both
bandits with fixed mutation rates, with and without ILP guidance. A
second ablation on the same pools isolates the ILP guidance: it
compares NSGA-III alone, the ILP solution as a seed only, the seed with
repair priority, and the full hybrid with the archive (UCB1, exact
enumeration disabled).

For out-of-distribution validation we use PrefLib
dataset~00036~\cite{mattei2013preflib}. Half of each instance's
vertices form the initial pool and the rest arrive over the $T$
rounds, and $n^*$ is scaled to $90/75$ of the initial pool. We run the
ILPs on all 179 instances with up to 294 vertices, and all
configurations, including both two-stage ILPs, on 12 instances of
128--147 vertices with 10 runs each at $\tau = 12$.

%% ============================================================
\section{Results and Discussion}
\label{sec:results}
%% ============================================================

Table~\ref{tab:main} summarizes the main experiment. All comparisons
are paired Wilcoxon tests over 100 runs with Bonferroni correction
unless stated otherwise. In 48\% of cleared rounds the feasible set
was small enough to enumerate exactly, the chain cap bound in 9.6\%,
and no implemented matching fell below the 90\% floor.

\begin{table*}[!t]
\centering
\caption{Trajectory results over 25 rounds (mean $\pm$ sd over 100
paired runs). $\Delta$: transplants lost against the myopic ILP.
High-cPRA lost: highly sensitized patients who withdrew.
D-/S-: dynamic/static NSGA.}
\label{tab:main}
\footnotesize
\setlength{\tabcolsep}{3.5pt}
\begin{tabular}{@{}cl cccccccc@{}}
\toprule
$\tau$ & Configuration & PoF (\%) & $\Delta$ (\%) & Transplants
 & High-cPRA tx & $F_2$ & $F_3$ & $F_4$ & High-cPRA lost \\
\midrule
\multirow{10}{*}{6}
 & S-NSGA-II + UCB1   & 1.63 $\pm$ 0.68 & $-$0.03 $\pm$ 1.74 & 173.6 $\pm$ 15.5 & 49.5 $\pm$ 6.9 & 0.354 & 0.060 & 0.097 & 25.6 $\pm$ 4.3 \\
 & S-NSGA-III fixed   & 1.56 $\pm$ 0.67 & $-$0.14 $\pm$ 1.50 & 173.8 $\pm$ 15.5 & 49.6 $\pm$ 7.0 & 0.356 & 0.061 & 0.099 & 25.6 $\pm$ 4.5 \\
 & D-UCB1 myopic      & 1.60 $\pm$ 0.61 & $-$0.06 $\pm$ 1.55 & 173.7 $\pm$ 15.1 & 49.5 $\pm$ 6.9 & 0.355 & 0.060 & 0.098 & 25.5 $\pm$ 4.4 \\
 & D-LinUCB myopic    & 1.59 $\pm$ 0.63 & $-$0.05 $\pm$ 1.74 & 173.7 $\pm$ 15.4 & 49.4 $\pm$ 6.9 & 0.355 & 0.060 & 0.097 & 25.7 $\pm$ 4.3 \\
 & D-UCB1 batched     & --              & 100.0              & 0.0              & 0.0            & --    & --    & --    & 65.6 $\pm$ 7.8 \\
 & D-LinUCB batched   & --              & 100.0              & 0.0              & 0.0            & --    & --    & --    & 65.6 $\pm$ 7.8 \\
 & Myopic ILP         & 0.00            & 0.00               & 173.6 $\pm$ 15.4 & 48.8 $\pm$ 7.1 & 0.369 & 0.077 & 0.137 & 26.3 $\pm$ 4.7 \\
 & Dickerson lex.     & 1.76 $\pm$ 1.27 & 0.15 $\pm$ 1.49    & 173.3 $\pm$ 15.5 & 59.2 $\pm$ 8.0 & 0.368 & 0.095 & 0.135 & 17.5 $\pm$ 4.5 \\
 & Two-stage (sum)    & 0.00            & $-$0.82 $\pm$ 1.48 & 175.0 $\pm$ 15.6 & 51.2 $\pm$ 7.5 & 0.356 & 0.058 & 0.087 & 24.2 $\pm$ 4.2 \\
 & Two-stage ($F_3$)  & 0.00            & $-$0.14 $\pm$ 1.49 & 173.8 $\pm$ 15.3 & 51.4 $\pm$ 7.3 & 0.369 & 0.057 & 0.096 & 24.0 $\pm$ 4.2 \\
\midrule
\multirow{10}{*}{12}
 & S-NSGA-II + UCB1   & 1.86 $\pm$ 0.71 & 0.17 $\pm$ 1.31   & 181.4 $\pm$ 15.9 & 54.4 $\pm$ 7.7 & 0.378 & 0.038 & 0.064 & 13.7 $\pm$ 2.8 \\
 & S-NSGA-III fixed   & 1.76 $\pm$ 0.76 & 0.16 $\pm$ 1.27   & 181.4 $\pm$ 15.9 & 54.3 $\pm$ 7.4 & 0.379 & 0.040 & 0.065 & 13.7 $\pm$ 3.1 \\
 & D-UCB1 myopic      & 1.81 $\pm$ 0.75 & 0.16 $\pm$ 1.26   & 181.4 $\pm$ 15.8 & 54.5 $\pm$ 7.9 & 0.379 & 0.039 & 0.061 & 13.6 $\pm$ 3.2 \\
 & D-LinUCB myopic    & 1.83 $\pm$ 0.71 & 0.15 $\pm$ 1.33   & 181.5 $\pm$ 15.9 & 54.2 $\pm$ 7.5 & 0.378 & 0.039 & 0.063 & 13.8 $\pm$ 3.0 \\
 & D-UCB1 batched     & 8.36 $\pm$ 1.02 & 24.39 $\pm$ 14.58 & 138.0 $\pm$ 30.9 & 41.6 $\pm$ 11.0 & 0.349 & 0.033 & 0.057 & 18.1 $\pm$ 4.4 \\
 & D-LinUCB batched   & 8.37 $\pm$ 1.01 & 23.32 $\pm$ 14.58 & 140.0 $\pm$ 31.3 & 42.2 $\pm$ 11.0 & 0.348 & 0.033 & 0.050 & 17.9 $\pm$ 4.4 \\
 & Myopic ILP         & 0.00            & 0.00              & 181.7 $\pm$ 15.7 & 53.9 $\pm$ 7.7 & 0.390 & 0.051 & 0.092 & 14.4 $\pm$ 3.3 \\
 & Dickerson lex.     & 2.12 $\pm$ 1.38 & 0.51 $\pm$ 1.55   & 180.8 $\pm$ 15.5 & 65.1 $\pm$ 8.0 & 0.390 & 0.087 & 0.091 & 7.5 $\pm$ 3.0 \\
 & Two-stage (sum)    & 0.00            & $-$0.37 $\pm$ 1.21 & 182.4 $\pm$ 15.9 & 55.7 $\pm$ 7.7 & 0.377 & 0.038 & 0.054 & 12.8 $\pm$ 3.3 \\
 & Two-stage ($F_3$)  & 0.00            & 0.08 $\pm$ 1.13   & 181.6 $\pm$ 15.7 & 56.0 $\pm$ 7.7 & 0.390 & 0.037 & 0.060 & 13.0 $\pm$ 3.2 \\
\midrule
\multirow{10}{*}{18}
 & S-NSGA-II + UCB1   & 1.91 $\pm$ 0.70 & 0.31 $\pm$ 1.31   & 183.3 $\pm$ 15.8 & 56.1 $\pm$ 8.0 & 0.388 & 0.034 & 0.056 & 7.1 $\pm$ 2.4 \\
 & S-NSGA-III fixed   & 1.80 $\pm$ 0.71 & 0.37 $\pm$ 1.27   & 183.2 $\pm$ 15.5 & 55.7 $\pm$ 7.5 & 0.390 & 0.035 & 0.057 & 7.2 $\pm$ 2.4 \\
 & D-UCB1 myopic      & 1.91 $\pm$ 0.71 & 0.27 $\pm$ 1.30   & 183.4 $\pm$ 15.8 & 56.0 $\pm$ 8.0 & 0.389 & 0.034 & 0.054 & 7.1 $\pm$ 2.5 \\
 & D-LinUCB myopic    & 1.90 $\pm$ 0.75 & 0.41 $\pm$ 1.34   & 183.2 $\pm$ 15.8 & 55.8 $\pm$ 7.7 & 0.389 & 0.034 & 0.055 & 7.2 $\pm$ 2.4 \\
 & D-UCB1 batched     & 8.10 $\pm$ 0.94 & 12.77 $\pm$ 7.73  & 160.4 $\pm$ 19.0 & 49.1 $\pm$ 7.7 & 0.389 & 0.033 & 0.053 & 8.1 $\pm$ 2.6 \\
 & D-LinUCB batched   & 8.25 $\pm$ 0.88 & 12.08 $\pm$ 6.92  & 161.7 $\pm$ 18.9 & 49.8 $\pm$ 8.1 & 0.388 & 0.033 & 0.049 & 8.0 $\pm$ 2.5 \\
 & Myopic ILP         & 0.00            & 0.00              & 183.9 $\pm$ 15.8 & 55.6 $\pm$ 8.0 & 0.400 & 0.045 & 0.080 & 7.9 $\pm$ 2.7 \\
 & Dickerson lex.     & 2.41 $\pm$ 1.66 & 0.66 $\pm$ 1.70   & 182.7 $\pm$ 15.6 & 67.0 $\pm$ 8.2 & 0.399 & 0.088 & 0.079 & 3.1 $\pm$ 2.0 \\
 & Two-stage (sum)    & 0.00            & $-$0.28 $\pm$ 1.23 & 184.4 $\pm$ 15.8 & 57.1 $\pm$ 7.8 & 0.386 & 0.033 & 0.048 & 6.8 $\pm$ 2.6 \\
 & Two-stage ($F_3$)  & 0.00            & 0.00 $\pm$ 1.30   & 183.9 $\pm$ 15.6 & 57.0 $\pm$ 7.6 & 0.399 & 0.032 & 0.050 & 7.1 $\pm$ 2.5 \\
\bottomrule
\end{tabular}
\end{table*}

%% ------------------------------------------------------------
\subsection{RQ1: Evolution of the Price of Fairness}
\label{sec:results:rq1}

Fig.~\ref{fig:pof} shows $\pof$ round by round. With myopic clearing,
the framework's mean $\pof$ was 1.60\%, 1.81\% and 1.91\% for
$\tau = 6$, 12 and 18. The first round, which clears the initial
backlog of 75 vertices, cost 8.44\%; after that the mean fell to
1.28--1.61\%. Over all rounds the per-run slope was negative
($p < 0.001$), but only because of the first round: over rounds 2--25
it did not differ from zero ($p = 0.42$--0.54). Once the backlog is
cleared, $\pof$ is stable over time.

What drives $\pof$ is the size of each round, not the passage of
time. Under the 90\% floor, $\pof$ can be positive only when the ILP
optimum reaches 10 transplants, which happened in 21--24\% of rounds.
In those rounds the framework gave up 7.7--7.8\% of transplants; in
all others it implemented a maximum matching.

The per-round price is also not a real loss over time. Summed over
the 25 rounds, the framework gave up 7.0, 7.7 and 8.0 transplants per
run against each round's optimum, yet it finished only $-0.06$, 0.31
and 0.50 transplants behind the myopic ILP. In other words,
94--100\% of the shortfall against each round's optimum was recovered
in later rounds, and $\Delta$ against the myopic ILP was not
significant at any $\tau$: patients who are not matched in one round
usually remain in the pool for the next.

\begin{figure*}[!htbp]
\centering
\includegraphics[width=\textwidth]{fig1_pof_per_round.png}
\caption{$\pof$ per cleared round (mean and 95\% CI over 100 runs) for
the framework (UCB1, LinUCB), the static NSGA-II baseline and the
lexicographic rule. After the first round, $\pof$ shows no trend.}
\label{fig:pof}
\end{figure*}

%% ------------------------------------------------------------
\subsection{RQ2: Myopic versus Batched Clearing}
\label{sec:results:rq2}

Batched clearing with $n^* = 90$ was worse at every $\tau$
(Fig.~\ref{fig:clearing}). At $\tau = 6$ the pool never reached 90
patients, so no one was transplanted and 65.6 highly sensitized
patients per run withdrew, against 25.5 under myopic clearing. At
$\tau = 12$ it transplanted 24\% fewer patients (138.0 vs.\ 181.4,
$d_z = 1.72$) and 24\% fewer highly sensitized patients, and 33\% more
of them withdrew. At $\tau = 18$ both losses were 12\% (all
$p_{\mathrm{bonf}} < 0.001$). Each batch was large, so $\pof$ sat near
the floor in every cleared round (8.1--8.4\%), while $F_3$ and $F_4$
were similar to myopic clearing. Waiting did not buy fairness; it only
cost transplants.

Better-tuned batching narrowed the gap without closing it
(Table~\ref{tab:batch}). Forcing a clearing at least every three
rounds avoided the failure at $\tau = 6$ and cut the loss to about
4.5\% at $\tau = 12$ and 18. Lowering the threshold to $n^* = 60$
reduced it to 5.5\% and 1.5\%. No variant transplanted more patients
than myopic clearing. This agrees with the near-optimality of greedy
matching without departure information~\cite{anderson2017efficient,
akbarpour2020thickness}, and, in this model, extends to the fairness
objectives.

\begin{figure*}[!htbp]
\centering
\includegraphics[width=\textwidth]{fig2_myopic_vs_batched.png}
\caption{Myopic (solid) versus batched (dashed) clearing:
$\pof$ per cleared round (top) and cumulative transplants (bottom),
mean over 100 runs. The UCB1 and LinUCB batched curves overlap. At
$\tau = 6$ batched clearing never clears.}
\label{fig:clearing}
\end{figure*}

\begin{table}[!t]
\centering
\caption{Batching variants (D-UCB1, 30 paired runs). Change in
transplants against myopic clearing on the same runs.}
\label{tab:batch}
\footnotesize
\setlength{\tabcolsep}{3pt}
\begin{tabular}{@{}lcccc@{}}
\toprule
Variant & $\tau$ & Transplants & High-cPRA tx & PoF (\%) \\
\midrule
Myopic                  & 12 & 178.6           & 54.9 & 1.67 \\
$n^* = 90$              & 12 & 126.1 ($-$29\%) & 38.9 & 8.55 \\
$n^* = 90$, max.\ 3     & 12 & 170.8 ($-$4.4\%) & 52.7 & 6.80 \\
$n^* = 60$              & 12 & 168.8 ($-$5.5\%) & 52.4 & 6.60 \\
\midrule
Myopic                  & 18 & 180.6           & 56.5 & 1.83 \\
$n^* = 90$              & 18 & 153.6 ($-$15\%) & 48.8 & 8.34 \\
$n^* = 90$, max.\ 3     & 18 & 172.4 ($-$4.5\%) & 54.3 & 6.85 \\
$n^* = 60$              & 18 & 177.9 ($-$1.5\%) & 56.2 & 4.32 \\
\bottomrule
\end{tabular}
\end{table}

%% ------------------------------------------------------------
\subsection{RQ3: Comparison with the Baselines}
\label{sec:results:rq3}

\emph{Myopic ILP.} The framework transplanted the same number of
patients at every $\tau$ ($p_{\mathrm{bonf}} = 1.00$) and was fairer
on all three fairness objectives ($p_{\mathrm{bonf}} < 0.001$;
Fig.~\ref{fig:f3}). It reduced $F_3$ by 21--24\%
($|d_z| = 1.4$--1.7), $F_4$ by 28--34\% ($|d_z| = 1.3$--1.8) and
$F_2$ by 3--4\% ($|d_z| = 1.9$--2.7), and 3--10\% fewer highly
sensitized patients withdrew ($p_{\mathrm{bonf}} \leq 0.032$). Part of
this gain costs nothing: with a 100\% floor ($\pof = 0$ in every
round) the framework still matched the ILP's volume (179.2 vs.\
178.4) and reduced $F_3$ by 18\% ($p_{\mathrm{bonf}} < 0.001$). The
volume ILP leaves fairness unused among the maximum matchings; the
two-stage ILP below uses it exactly.

\emph{Two-stage ILP.} The exact two-stage ILP was more effective than
the framework at maximum volume, at all three levels of comparison
(Table~\ref{tab:levels}, Fig.~\ref{fig:levels}).
\begin{itemize}
  \item \emph{Identical pools.} The two-stage ILP (sum) reached
    $F_3 = 0.0217$ and $F_4 = 0.0348$, against 0.0548 and 0.1143 for
    the framework's maximum-volume choice (both
    $p_{\mathrm{bonf}} \leq 0.006$). The framework's front held only
    1.1 solutions at maximum volume per pool on average, and its
    lowest $F_3$ at maximum volume equalled the optimum of the
    two-stage ILP ($F_3$) in half of the pools.
  \item \emph{Fixed trajectories.} On the same pools in every round
    ($\tau = 12$, 689 rounds), the gap was smaller. On the stage-2
    score~\eqref{eq:twostage}, the volume ILP scored 3.000, the
    two-stage ILP 2.322 and the framework's maximum-volume choice
    2.606. The framework's lowest $F_3$ at maximum volume equalled the
    optimum of the two-stage ILP ($F_3$) in 95.4\% of rounds (100\% of
    enumerated rounds), so
    most of the remaining gap comes from the selection rule, which
    differs from the stage-2 objective, rather than from the search.
  \item \emph{Full trajectories.} Against the framework with its 90\%
    floor, the two-stage ILP made 1.3, 1.0 and 1.0 more transplants
    per run ($\tau = 6$, 12, 18; all $p_{\mathrm{bonf}} < 0.001$) and
    1.1--1.7 more highly sensitized transplants
    ($p_{\mathrm{bonf}} < 0.001$). Its $F_4$ was lower at every $\tau$
    ($p_{\mathrm{bonf}} < 0.001$), its $F_2$ was lower at $\tau = 12$
    and 18 ($p_{\mathrm{bonf}} \leq 0.040$), and its $F_3$ was lower at
    $\tau = 6$ ($p_{\mathrm{bonf}} = 0.034$) and did not differ at
    $\tau = 12$ and 18. With a 100\% floor, the framework matched the
    two-stage ILP's volume but had higher $F_2$, $F_3$ and $F_4$
    ($\tau = 12$, 30 runs, $p_{\mathrm{bonf}} \leq 0.004$).
\end{itemize}
The two-stage ILP is therefore the stronger method when the aim is
the fairest matching at maximum volume. The 90\% floor gives up up to
10\% of a round's transplants, but over full trajectories this did not
make the framework fairer than the two-stage ILP. Against the myopic
ILP, the two-stage ILP was much fairer and made as many or more
transplants (more at $\tau = 6$, $p_{\mathrm{bonf}} < 0.001$).

\emph{Tie-breaking.} The deterministic HLA tie-break rarely affects the
reported outcomes. In 956 enumerated rounds with an exact front,
24\% had more than one maximum matching with the optimal stage-2
score, but the tied matchings differed in $F_2$, $F_3$ or $F_4$ in only
0.1\% of rounds (mean range of $F_3$ among them: $5 \times 10^{-5}$).

\emph{Lexicographic rule}~\cite{dickerson2014price}. Volume and
$\pof$ did not differ from the framework. The lexicographic rule
transplanted about 20\% more highly sensitized patients and lost far
fewer of them to withdrawal (3.1--17.5 vs.\ 7.1--25.5 per run), but it
over-corrected: their match rate rose well above everyone else's, so
its $F_3$ was the largest of all myopic methods (0.087--0.095), and
$F_2$ and $F_4$ did not improve on the ILP. The two methods pursue
different goals, maximum priority versus balanced access, and the
choice between them is a policy decision.

\emph{Static baselines.} Static NSGA-II with UCB1 and static NSGA-III
with a fixed rate gave the same volume, $\pof$ and fairness as the
dynamic framework. Warm-starting added only 0.55\% hypervolume at
$\tau = 18$ ($p_{\mathrm{bonf}} = 0.007$). With ILP guidance and exact
small rounds, little is left for the warm start to improve.

\begin{table}[!t]
\centering
\caption{Framework and ILP baselines at three levels of comparison
(mean $\pm$ sd). A: same pool, one round (30 pools). B: same pool in
every round of the framework's trajectories ($\tau = 12$, per-run
means). C: each method clears its own pools ($\tau = 12$, per-run
means; framework 90\% floor and ILPs over 100 runs, 100\% floor over
30). Volume ILP, two-stage ILP and the 100\% pick or floor are
restricted to maximum volume; the 90\% pick or floor may give up to
10\% of a round's transplants. Score: stage-2
objective~\eqref{eq:twostage} without the tie-break.}
\label{tab:levels}
\footnotesize
\setlength{\tabcolsep}{2.5pt}
\begin{tabular}{@{}lcccc@{}}
\toprule
Method & $F_2$ & $F_3$ & $F_4$ & Score \\
\midrule
\multicolumn{5}{@{}l}{\textit{A. Identical pools}} \\
Volume ILP          & 0.350 & 0.056 & 0.126 & 2.875 \\
Two-stage ILP (sum) & 0.324 & 0.022 & 0.035 & 1.826 \\
Two-stage ILP ($F_3$) & 0.347 & 0.022 & 0.047 & 1.977 \\
Framework, 100\% pick & 0.348 & 0.055 & 0.114 & 2.783 \\
Framework, 90\% pick  & 0.330 & 0.023 & 0.055 & 2.212 \\
\multicolumn{5}{@{}l}{\textit{B. Fixed trajectories}} \\
Volume ILP          & 0.382 & 0.049 & 0.090 & 2.753 \\
Two-stage ILP (sum) & 0.378 & 0.036 & 0.054 & 2.356 \\
Two-stage ILP ($F_3$) & 0.382 & 0.035 & 0.058 & 2.414 \\
Framework, 100\% pick & 0.379 & 0.041 & 0.069 & 2.681 \\
Framework, 90\% pick  & 0.378 & 0.037 & 0.061 & 2.578 \\
\multicolumn{5}{@{}l}{\textit{C. Full trajectories}} \\
Myopic ILP          & 0.390 & 0.051 & 0.092 & -- \\
Two-stage ILP (sum) & 0.377 & 0.038 & 0.054 & -- \\
Two-stage ILP ($F_3$) & 0.390 & 0.037 & 0.060 & -- \\
Framework, 100\% floor & 0.381 & 0.041 & 0.069 & -- \\
Framework, 90\% floor  & 0.379 & 0.039 & 0.061 & -- \\
\bottomrule
\multicolumn{5}{@{}l}{Score in B: mean of per-run means over all rounds.}
\end{tabular}
\end{table}

\begin{figure*}[!htbp]
\centering
\includegraphics[width=\textwidth]{fig3_three_levels.png}
\caption{$F_2$, $F_3$ and $F_4$ of the framework and the ILP baselines
at the three levels of comparison (mean and 95\% CI).}
\label{fig:levels}
\end{figure*}

\begin{figure*}[!htbp]
\centering
\includegraphics[width=\textwidth]{fig6_forest_tau12.png}
\caption{Paired differences, framework (UCB1, 90\% floor) minus each
baseline, over 100 runs at $\tau = 12$ (mean and 95\% CI). Negative
values favor the framework. Filled markers: significant after
Bonferroni correction.}
\label{fig:forest}
\end{figure*}

\begin{figure*}[!htbp]
\centering
\includegraphics[width=\textwidth]{trajectory_F3.png}
\caption{High-cPRA access gap $F_3$ per cleared round (mean and 95\%
CI over 100 runs). The framework stays below the myopic ILP; the
lexicographic rule over-corrects. Batched bands are wider because few
runs clear in some rounds.}
\label{fig:f3}
\end{figure*}

%% ------------------------------------------------------------
\subsection{RQ4: Trade-off Solutions below Maximum Volume}
\label{sec:results:rq4}

A single two-stage ILP solution fixes the volume at its maximum. The
front, in contrast, contains matchings at several volumes. We asked
whether these give meaningful options and how close they are to the
exact answer at the same volume (Table~\ref{tab:tradeoff},
Fig.~\ref{fig:tradeoff}).

\emph{Identical pools.} The number of front solutions grew from 3.4 at
the 95\% level to 23.0 at 80\%. In 23\%, 53\%, 80\% and 90\% of pools
(95\% to 80\% levels), the front held a solution with a lower stage-2
score than the two-stage ILP at maximum volume, and in 3\%, 13\%, 33\%
and 60\% of pools a solution with lower or equal $F_2$, $F_3$ and
$F_4$ at once. The exact two-stage ILP at the same volumes was still
fairer at every level (score 1.291 against 1.318 at 80\%; 1.414
against 1.790 at 90\%), although 97--100\% of these solves stopped at
the 120~s limit. Each exact solve gives one matching at one volume;
the front gave all of them in one run of 27.6~s.

\emph{Fixed trajectories.} On the pools of real trajectory rounds the
front was close to the exact reference at every level. At $\tau = 12$
the fairest front solution scored 2.416 and 2.167 at the 90\% and
80\% levels, against 2.379 and 2.146 for the exact reference, and its
$F_3$ was the same (0.035 at 90\%; 0.023 against 0.024 at 80\%).
Where the level allowed fewer transplants than $F_1^{*}$ (70\% of
rounds at the 80\% level), the front offered 7.5 options on average,
and in 64\% of these rounds at least one was fairer than the two-stage
ILP at maximum volume. At the 80\% level, the fairest front solution
had significantly lower $F_2$, $F_3$ and $F_4$ than the two-stage ILP
at maximum volume at every $\tau$ ($p_{\mathrm{bonf}} < 0.001$). The
floor rule, however, did not pick these solutions: its choice scored
2.540 at 80\% and 2.830 at maximum volume, against 2.167 and 2.548 for
the fairest solution on the same front. The front holds near-exact
trade-off solutions; the selection rule loses most of the fairness.

\emph{Selection rules over full trajectories.} The rules move the
framework along a consistent trade-off at all three $\tau$
(Table~\ref{tab:select}, Fig.~\ref{fig:select}). The 100\% floor matched
the two-stage ILP's volume but was less fair. The 80\% floor gave a
lower $F_3$ than the two-stage ILP at every $\tau$ (0.045, 0.028 and
0.025 against 0.054, 0.035 and 0.032; $p_{\mathrm{bonf}} < 0.001$) and
a lower $F_2$ and $F_4$ at $\tau = 6$ and 12, for 2.1--2.9\% fewer
transplants ($\Delta = 1.7$--2.1\%). The compromise rule reduced $F_3$
to 0.013--0.026 at a much larger cost ($\pof = 36$--37\%,
$\Delta = 5.7$--6.5\%). Both rules also transplanted fewer highly
sensitized patients than the two-stage ILP (2.2--2.6 fewer with the
80\% floor), because $F_3$ can fall when fewer patients are matched
overall. These options are not available from a single two-stage ILP
solution; whether a given trade-off is acceptable is a policy choice.

\begin{table*}[!t]
\centering
\caption{Fairness available below maximum volume. Level $\rho$ allows
$V = \lceil \rho F_1^{*} \rceil$ or more transplants. Options: distinct
front solutions with $F_1 \geq V$. Score: stage-2
objective~\eqref{eq:twostage} (lower = fairer). Exact: best stage-2
solution among the volumes solved from $V$ to $F_1^{*}$ (one solve per
level, time-limited). Means over all units.}
\label{tab:tradeoff}
\footnotesize
\setlength{\tabcolsep}{3pt}
\begin{tabular}{@{}lccccccc@{}}
\toprule
Level $\rho$ & Options & Front best & Rule pick & Exact & Front best $F_3$ & Exact $F_3$ & Time-limit share \\
\midrule
\multicolumn{8}{@{}l}{\textit{Identical pools (30 pools)}} \\
100\% & 1.1 & 2.762 & 2.783 & 1.826 & 0.055 & 0.022 & 0.77 \\
95\% & 3.4 & 2.216 & 2.509 & 1.519 & 0.034 & 0.010 & 0.97 \\
90\% & 8.5 & 1.790 & 2.212 & 1.414 & 0.015 & 0.006 & 0.97 \\
85\% & 15.7 & 1.459 & 1.812 & 1.347 & 0.008 & 0.005 & 1.00 \\
80\% & 23.0 & 1.318 & 1.695 & 1.291 & 0.003 & 0.004 & 1.00 \\
\multicolumn{8}{@{}l}{\textit{Fixed trajectories, $\tau = 12$ (750 rounds)}} \\
100\% & 2.0 & 2.548 & 2.830 & 2.476 & 0.041 & 0.039 & 0.05 \\
95\% & 2.2 & 2.523 & 2.819 & 2.460 & 0.040 & 0.038 & 0.05 \\
90\% & 3.0 & 2.416 & 2.717 & 2.379 & 0.035 & 0.035 & 0.08 \\
85\% & 4.2 & 2.307 & 2.634 & 2.281 & 0.030 & 0.030 & 0.09 \\
80\% & 5.9 & 2.167 & 2.540 & 2.146 & 0.023 & 0.024 & 0.12 \\
\bottomrule
\end{tabular}
\end{table*}

\begin{figure*}[!htbp]
\centering
\includegraphics[width=\textwidth]{fig4_tradeoff_levels.png}
\caption{Fairness against the volume level $\rho$: fairest front
solution with $F_1 \geq V$, the floor-rule choice, the exact two-stage
ILP (best of the volumes solved from $V$ to $F_1^{*}$) and the
two-stage ILP at maximum volume (dotted). Mean and 95\% CI.}
\label{fig:tradeoff}
\end{figure*}

\begin{table*}[!t]
\centering
\caption{Selection rules over full trajectories compared with the
two-stage ILP and the myopic ILP (mean $\pm$ sd over the same 30 paired
runs).}
\label{tab:select}
\footnotesize
\setlength{\tabcolsep}{3pt}
\begin{tabular}{@{}lccccccc@{}}
\toprule
Configuration & Tx & High-cPRA tx & $F_2$ & $F_3$ & $F_4$ & PoF (\%) & $\Delta$ (\%) \\
\midrule
\multicolumn{8}{@{}l}{\textit{$\tau = 6$}} \\
Framework, 100\% floor & 171.7 & 50.6 & 0.359 & 0.065 & 0.106 & 0.00 & $-$1.02 \\
Framework, 90\% floor  & 170.2 & 49.9 & 0.355 & 0.060 & 0.101 & 1.46 & $-$0.13 \\
Framework, 80\% floor  & 166.9 & 49.3 & 0.345 & 0.045 & 0.077 & 9.86 & 1.81 \\
Framework, compromise  & 158.9 & 46.6 & 0.320 & 0.026 & 0.047 & 36.04 & 6.55 \\
Two-stage ILP (sum)    & 171.9 & 51.5 & 0.355 & 0.054 & 0.091 & 0.00 & $-$1.12 \\
Myopic ILP             & 170.0 & 48.9 & 0.370 & 0.075 & 0.140 & 0.00 & 0.00 \\
\multicolumn{8}{@{}l}{\textit{$\tau = 12$}} \\
Framework, 100\% floor & 179.2 & 54.9 & 0.381 & 0.041 & 0.069 & 0.00 & $-$0.49 \\
Framework, 90\% floor  & 178.6 & 54.9 & 0.378 & 0.037 & 0.061 & 1.68 & $-$0.10 \\
Framework, 80\% floor  & 175.3 & 53.3 & 0.371 & 0.028 & 0.048 & 10.34 & 1.74 \\
Framework, compromise  & 168.3 & 51.1 & 0.357 & 0.015 & 0.028 & 37.09 & 5.70 \\
Two-stage ILP (sum)    & 179.1 & 55.7 & 0.375 & 0.035 & 0.056 & 0.00 & $-$0.43 \\
Myopic ILP             & 178.4 & 53.9 & 0.389 & 0.050 & 0.093 & 0.00 & 0.00 \\
\multicolumn{8}{@{}l}{\textit{$\tau = 18$}} \\
Framework, 100\% floor & 181.8 & 56.7 & 0.392 & 0.037 & 0.063 & 0.00 & $-$0.72 \\
Framework, 90\% floor  & 180.6 & 56.5 & 0.388 & 0.034 & 0.058 & 1.83 & $-$0.05 \\
Framework, 80\% floor  & 176.8 & 54.9 & 0.382 & 0.025 & 0.043 & 10.62 & 2.06 \\
Framework, compromise  & 169.9 & 52.3 & 0.370 & 0.013 & 0.025 & 37.43 & 5.92 \\
Two-stage ILP (sum)    & 181.5 & 57.5 & 0.385 & 0.032 & 0.048 & 0.00 & $-$0.53 \\
Myopic ILP             & 180.5 & 55.5 & 0.400 & 0.045 & 0.081 & 0.00 & 0.00 \\
\bottomrule
\end{tabular}
\end{table*}

\begin{figure*}[!htbp]
\centering
\includegraphics[width=\textwidth]{fig5_selection_rules.png}
\caption{Selection rules over full trajectories: trajectory loss
$\Delta$ against $F_3$ (top) and $F_4$ (bottom), mean and 95\% CI over
the same 30 paired runs. The framework's rules (blue, joined) run from
the 100\% floor to the compromise rule.}
\label{fig:select}
\end{figure*}

%% ------------------------------------------------------------
\subsection{RQ5: Operator Selection and ILP Guidance}
\label{sec:results:rq5}

UCB1 and LinUCB were statistically identical on every outcome
($p_{\mathrm{bonf}} = 1.00$). UCB1 spread its choices almost evenly
over the four mutation rates (normalized entropy 1.00), whereas LinUCB
leaned slightly towards $p_m = 0.25$ under myopic clearing (28--29\%)
and towards $p_m = 0.01$ under batched clearing (33--38\%)
(Fig.~\ref{fig:arms}).
Table~\ref{tab:hybrid} isolates the parts of the ILP guidance. NSGA-III
alone never reached the volume optimum, and its selected solutions had
a mean $\pof$ of 28.3\%, against 8.2\% for the full hybrid
($p_{\mathrm{bonf}} < 0.001$); roughly 20 percentage points of the
unguided $\pof$ were search failure rather than fairness. The ILP
solution as a seed alone changed nothing ($p_{\mathrm{bonf}} = 1.00$
for every metric), because the greedy repair rebuilt it into a smaller
matching. Repair priority was the decisive part, raising optimum
attainment to 77\% of pools, and the archive closed the remaining gap
to 100\% with a slightly higher hypervolume ($p_{\mathrm{bonf}} =
0.007$). The guidance cost 5.4~s per round over NSGA-III alone in
these pools, 2.4~s of it for the ILP. In the main experiment the ILP
took 0.48~s per round, 2.1\% of a framework clearing (23.3~s), and
the HiGHS fallback was never needed.

\begin{table}[!t]
\centering
\caption{Ablation of the ILP guidance: one round on 30 pools of 75
vertices (UCB1, $N = 100$, $G_{\max} = 150$). ILP optimum: 42.3
transplants. PoF of the solution chosen by the 90\% floor rule,
relative to the ILP optimum.}
\label{tab:hybrid}
\footnotesize
\setlength{\tabcolsep}{2.5pt}
\begin{tabular}{@{}lccccc@{}}
\toprule
Arm & At optimum & Max $F_1$ & HV & PoF (\%) & Time (s) \\
\midrule
NSGA-III alone   & 0 / 30  & 32.6 & 0.566 & 28.3 & 18.2 \\
ILP seed only    & 0 / 30  & 32.7 & 0.566 & 28.8 & 20.9 \\
Seed + priority  & 23 / 30 & 41.8 & 0.696 & 9.2  & 25.1 \\
Full hybrid      & 30 / 30 & 42.3 & 0.703 & 8.2  & 23.6 \\
\bottomrule
\multicolumn{6}{@{}l}{Time: MOEA plus ILP (2.4~s) where the ILP is used.}
\end{tabular}
\end{table}
The ablation in Table~\ref{tab:ablation} explains why the choice of
bandit mattered so little. Without ILP guidance, no setting reached
the optimum in any of the 30 pools (32.7 vs.\ 42.3 transplants); with
guidance, every setting reached it in every pool. With guidance, UCB1
beat fixed rates of 0.10 and 0.25 but was slightly below a fixed rate
of 0.01, and the two bandits did not differ ($p_{\mathrm{bonf}} =
0.11$). All differences in hypervolume were below 3\%. The bandits
protect against a poorly chosen rate without tuning, but here they do
not beat the best fixed rate, and the pool context gives LinUCB no
advantage. The ILP guidance, not the operator selection, is what
drives search quality.

\begin{figure*}[!htbp]
\centering
\includegraphics[width=\textwidth]{aos_arm_selection.png}
\caption{Share of generations using each mutation rate per round
($\tau = 12$).}
\label{fig:arms}
\end{figure*}

\begin{table}[!t]
\centering
\caption{AOS ablation: one round on 30 pools of 75 vertices.
ILP optimum: 42.3 transplants.}
\label{tab:ablation}
\footnotesize
\setlength{\tabcolsep}{3pt}
\begin{tabular}{@{}lcccc@{}}
\toprule
& \multicolumn{2}{c}{Without guidance} & \multicolumn{2}{c}{With guidance} \\
\cmidrule(lr){2-3}\cmidrule(l){4-5}
Setting & HV & Max $F_1$ & HV & Max $F_1$ \\
\midrule
UCB1          & 0.568 & 32.7 & 0.703 & 42.3 \\
LinUCB        & 0.567 & 32.6 & 0.704 & 42.3 \\
Fixed $0.01$  & 0.570 & 32.7 & 0.708 & 42.3 \\
Fixed $0.10$  & 0.564 & 32.7 & 0.698 & 42.3 \\
Fixed $0.25$  & 0.558 & 32.6 & 0.688 & 42.3 \\
\midrule
Pools at optimum & \multicolumn{2}{c}{0 / 30} & \multicolumn{2}{c}{30 / 30} \\
\bottomrule
\end{tabular}
\end{table}

%% ------------------------------------------------------------
\subsection{RQ6: Sensitivity and Scalability}
\label{sec:results:rq6}

Table~\ref{tab:sens} summarizes the sensitivity analyses at
$\tau = 12$; the selection rules are discussed in
Section~\ref{sec:results:rq4}.

\emph{CIT, exit rate and arrival rate.} A 12~h CIT limit removed
13--14\% of transplants for every method, but the framework still
matched the ILP and reduced $F_3$. Changing $\mu_d$ left $\pof$
unchanged, although a higher exit rate hurt batched clearing far more
($-26\%$) than myopic clearing ($-1.6\%$). At $\lambda = 7$ and 13 the
framework again matched the ILP and reduced $F_3$ by 22\% and 27\%;
$\pof$ rose with pool thickness (1.03\%, 1.67\%, 3.04\%), in line with
RQ1.

\emph{Scalability.} With initial pools of 200 vertices, the framework
took 37~s per round, against 5~s for the ILP and 8~s for the
lexicographic rule. It kept the ILP's volume per round (27.6 vs.\
27.7) and reduced $F_3$ by 39\%; its mean $\pof$ rose to 7.0\%
because most rounds now exceeded 10 transplants. With 500 initial
vertices the ILP took 38~s per round, and the lexicographic rule's
over-correction grew ($F_3 = 0.46$ against 0.05 for the ILP). The
framework was not run at 500 vertices because of memory limits.

\begin{table}[!t]
\centering
\caption{Sensitivity at $\tau = 12$ (30 paired runs): D-UCB1 myopic
against the myopic ILP.}
\label{tab:sens}
\footnotesize
\setlength{\tabcolsep}{2.5pt}
\begin{tabular}{@{}lccccc@{}}
\toprule
Variant & PoF (\%) & Tx (fw) & Tx (ILP) & $F_3$ (fw) & $F_3$ (ILP) \\
\midrule
Base                 & 1.67  & 178.6 & 178.4 & 0.037 & 0.050 \\
Floor 100\%          & 0.00  & 179.2 & 178.4 & 0.041 & 0.050 \\
Floor 80\%           & 10.34 & 175.3 & 178.4 & 0.029 & 0.050 \\
Compromise           & 37.09 & 168.3 & 178.4 & 0.015 & 0.050 \\
$\cit_{\max} = 16$~h & 1.78  & 178.0 & 177.8 & 0.036 & 0.049 \\
$\cit_{\max} = 12$~h & 1.51  & 154.6 & 153.9 & 0.033 & 0.040 \\
$\mu_d = 0.01$       & 1.75  & 179.4 & 179.1 & 0.036 & 0.048 \\
$\mu_d = 0.04$       & 1.69  & 175.7 & 176.0 & 0.040 & 0.054 \\
$\lambda = 7$        & 1.03  & 137.5 & 137.1 & 0.038 & 0.049 \\
$\lambda = 13$       & 3.04  & 224.7 & 225.2 & 0.036 & 0.049 \\
\bottomrule
\multicolumn{6}{@{}l}{fw: framework; Tx: transplants.}
\end{tabular}
\end{table}

%% ------------------------------------------------------------
\subsection{Computational Cost}
\label{sec:results:cost}

Table~\ref{tab:cost} reports the cost of every method under the same
conditions. On trajectory rounds the two-stage ILP (sum) needed
5.08~s per round at $\tau = 12$, against 19.30~s for the framework,
0.52~s for the volume ILP and 1.14~s for the lexicographic rule, and
its stage~2 reached the 60~s limit in 4.0\% of rounds. On the larger
75-vertex pools of the identical-pool comparison, stage~2 took
233.3~s on average and reached its 300~s limit in 67\% of pools,
against 27.2~s for the framework. Exact solves at lower volume levels
are more expensive still: one solve per level took 14--63~s on
trajectory rounds and about 125~s on identical pools, where almost all
reached the limit. The exact method is thus fast on the small pools
that arise once the backlog is cleared, but its cost grows quickly
with pool size and with the number of volume levels requested, while
the framework returns a whole front at a roughly constant cost.

\begin{table*}[!t]
\centering
\caption{Computational cost. Identical pool: one round on a 75-vertex
pool (stage-2 limit 300~s). Trajectory: mean seconds per cleared round
($\tau = 12$; stage-2 limit 60~s). One core per run.}
\label{tab:cost}
\footnotesize
\setlength{\tabcolsep}{3pt}
\begin{tabular}{@{}lcccc@{}}
\toprule
Method & Identical pool (s) & Time-limit share & Trajectory (s/round) & PrefLib (s/round) \\
\midrule
Volume (myopic) ILP & 3.2 & -- & 0.52 & 0.16 \\
Two-stage ILP (sum) & 233.3 & 0.67 & 5.08 & 2.45 \\
Two-stage ILP ($F_3$) & 58.5 & 0.17 & 1.60 & 0.76 \\
Framework (NSGA-III + AOS, incl.\ ILP) & 27.2 & -- & 19.30 & 4.41 \\
Dickerson lex. & -- & -- & 1.14 & 0.38 \\
Exact stage 2 at all lower levels & -- & 0.05--0.12\textsuperscript{a} & 24.96 & -- \\
\bottomrule
\multicolumn{5}{@{}l}{\textsuperscript{a}\,Range over the four lower levels at $\tau = 12$ (Table~\ref{tab:tradeoff}).}
\end{tabular}
\end{table*}

%% ------------------------------------------------------------
\subsection{Out-of-Distribution Validation on PrefLib}
\label{sec:results:ood}

Table~\ref{tab:ood} reports 120 paired runs. With myopic clearing the
framework transplanted 75.03 patients against 75.38 for the ILP, a
difference that was not significant ($p_{\mathrm{bonf}} = 0.19$), at a
mean $\pof$ of 0.73\%. Its $F_3$ was 4\% lower than the ILP's, but not
significantly; its $F_2$ was significantly lower
(Fig.~\ref{fig:forestood}). The two-stage ILP (sum) was again the
stronger method at maximum volume: it made more transplants (75.79),
had a lower $\Delta$ ($-0.50\%$ against 0.50\%) and had lower $F_2$
and $F_3$ than the framework (all significant after Bonferroni
correction). The lexicographic rule again transplanted the most highly
sensitized patients and had the largest gap among the myopic methods.
Batched clearing with a three-round maximum interval transplanted
3.9\% fewer patients than myopic clearing and doubled $F_3$ (both
$p_{\mathrm{bonf}} < 0.001$). On all 179 instances, the lexicographic
rule cost 0.84\% of transplants on average, rising to 2.2\% at 294
vertices.

The findings on volume and clearing policy carry over to PrefLib, but
the framework's fairness gain is much smaller. PrefLib graphs are
dense, so most patients are matched in the first round and few rounds
offer a choice between maximum matchings. Centers are assigned at
random, so $F_4$ is not reported.

\begin{table}[!t]
\centering
\caption{PrefLib 00036: 12 instances of 128--147 vertices, 10 paired
runs each, $\tau = 12$ (mean $\pm$ sd).}
\label{tab:ood}
\footnotesize
\setlength{\tabcolsep}{2.5pt}
\begin{tabular}{@{}lccccc@{}}
\toprule
Configuration & PoF (\%) & Transplants & High-cPRA tx & $F_2$ & $F_3$ \\
\midrule
Myopic ILP            & 0.00 & 75.38 $\pm$ 13.46 & 10.72 $\pm$ 2.63 & 0.325 & 0.062 \\
Dickerson lex.        & 1.54 & 75.09 $\pm$ 13.59 & 13.57 $\pm$ 2.96 & 0.324 & 0.080 \\
Two-stage (sum)       & 0.00 & 75.79 $\pm$ 13.81 & 11.40 $\pm$ 2.71 & 0.312 & 0.055 \\
Two-stage ($F_3$)     & 0.00 & 75.36 $\pm$ 13.66 & 11.45 $\pm$ 2.64 & 0.324 & 0.056 \\
D-UCB1 myopic         & 0.73 & 75.03 $\pm$ 13.68 & 10.59 $\pm$ 2.53 & 0.316 & 0.060 \\
D-LinUCB myopic       & 0.77 & 74.96 $\pm$ 13.69 & 10.62 $\pm$ 2.63 & 0.315 & 0.059 \\
D-UCB1 batched\textsuperscript{a} & 2.20 & 72.12 $\pm$ 13.28 & 10.16 $\pm$ 2.54 & 0.325 & 0.118 \\
\bottomrule
\multicolumn{6}{@{}l}{\textsuperscript{a}\,three-round maximum interval.}
\end{tabular}
\end{table}

\begin{figure*}[!htbp]
\centering
\includegraphics[width=\textwidth]{fig7_forest_preflib.png}
\caption{PrefLib 00036: paired differences, framework (UCB1, 90\%
floor) minus each baseline, over 120 paired units (mean and 95\% CI).
Filled markers: significant after Bonferroni correction.}
\label{fig:forestood}
\end{figure*}

%% ------------------------------------------------------------
\subsection{Limitations}
\label{sec:results:limits}

All results come from synthetic simulation informed by public ANZKX
information, not from validation on patient-level clinical data. The
exit rate, withdrawal thresholds, transport and CIT constraints and
the HLA model are modeling assumptions rather than validated ANZKX
parameters, so the results describe model behaviour, not clinical
outcomes. No single arrival rate matches both the ANZKX volume and
pool size, although the findings held at $\lambda = 7$, 10 and 13.
$F_4$ is a proxy. $\pof$ is bounded by the 90\% floor, so its level
partly reflects the selection rule, and the rule, which differs from
the stage-2 objective, did not choose the fairest solutions available
on the front. Batched clearing was tested without information about
upcoming departures. The two-stage ILP and the exact level solves were
time-limited, so on large pools their solutions are the best found,
not proven optima. The framework's fairness gain over the volume ILP
was small on PrefLib, and the framework was not run at 500 vertices.
The ILP guidance reduces search-failure bias in $\pof$, but the
remaining loss also depends on the selection rule and on how fully the
rest of the front is approximated.

%% ============================================================
\section{Conclusion}
\label{sec:conclusion}
%% ============================================================

This paper asked how the price of fairness behaves when a kidney
exchange is cleared month after month rather than once, whether
myopic or batched clearing better preserves the balance between
efficiency and equity, and what a multi-objective search adds to an
exact baseline. We studied this in an ANZKX-calibrated simulator with
withdrawals, deceased-donor exits and cold ischemia limits across six
regional nodes, using an ILP-guided NSGA-III with bandit-based
operator selection.

Four findings stand out. First, the per-round price of fairness was
small and did not accumulate. Myopic clearing with the framework
transplanted as many patients as the volume-maximizing ILP, while
reducing the high-cPRA access gap by 21--24\%, the logistics gap by
28--34\% and waiting-time inequality by 3--4\%. The price of fairness
averaged 1.6--1.9\% per round, stayed flat after the first round, and
94--100\% of the per-round shortfall was recovered in later rounds.
Second, batching did not help. With a threshold of 90 patients it
never cleared at $\tau = 6$ and lost 12--24\% of transplants at
$\tau = 12$ and 18, without improving fairness; better settings
narrowed the gap but never closed it. Third, an exact two-stage ILP
was the stronger method at maximum volume: on identical pools, on
fixed trajectories and over full trajectories it selected fairer
maximum matchings, and over full trajectories it also made 1.0--1.3
more transplants per run than the framework with its 90\% floor.
Fourth, the evolutionary front was useful below maximum volume. On
trajectory rounds it contained solutions close to the exact
reference at each volume level, and an 80\% floor lowered the access
gap by 17--23\% relative to the two-stage ILP for 2.1--2.9\% fewer
transplants at every withdrawal threshold. The floor rule did not
always choose the fairest of these solutions, and the ILP guidance,
not the bandit, drove the quality of the search.

Within this computational model, matching every month outperformed
waiting for a thicker pool. When the aim is the fairest matching at
maximum volume, an exact two-stage ILP is the better tool; when some
volume may be given up for fairness, the front provides the range of
options from which a policy can be chosen. These are model results;
whether they hold in practice must be established with program data
before any claim about clinical use.

Future work should test the framework on ANZDATA records, subject to
ethics approval, design selection rules aligned with the stage-2
objective, compare the front with exact $\varepsilon$-constraint
fronts, extend the methods to larger pools with column generation,
and study clearing policies that anticipate departures.

\balance
\bibliographystyle{IEEEtran}
\bibliography{Paper2_References}

\end{document}
